% SOURCE docs/release/2026-10-29/source/docs/writing-repair/2026-10-29/src/9ac4870b17074c5a.tex CHAR 0
\chapter{Longitudinal and survival software as a decision model}
\label{exec17:s5}

A bank comparing repeated servicing policies needs more than a curve with a
confidence ribbon. An account can default, repay early, cease to be observed,
or receive a different action after its financial condition changes. A policy
that reduces default incidence can also shorten the profitable relationship.
A policy that improves every baseline risk stratum can have a larger late
hazard among its remaining accounts. Two programs can return numerically
similar objects named survival while counting different accounts. These
differences affect both the scientific question and the financial decision.

The longitudinal identification and sequential-estimation results developed
earlier supply the foundation. Here the question is how to preserve that
foundation when choosing and combining software. The argument proceeds from
the event clock and supported horizon to survivor selection, competing-event
simulation, the meaning of a modified action, and uncertainty in the resulting
monetary contrast. Its finite laws and monetary amounts are stipulated,
reproducible constructions. They provide exact comparison references and
expose otherwise hidden distinctions; they are not bank observations or new
executions of the external estimators.~\citep{exec17_s5_math}

\section{The decision clock and the unsupported tail}
\label{exec17:s5-horizon}

Let $j=1,\ldots,K$ index servicing periods. At the opening of a period the
bank observes the permitted history $H_j$, chooses an action $A_j$, and then
observes transitions and exits during that period. An intervention must specify
what the bank would do using the history available at that opening. A balance
posted after the decision, a later repayment classification, and a corrected
record unavailable to the decision maker cannot enter that history merely
because all three eventually appear in a warehouse. The data schema must
represent availability as well as the time to which a value refers.

Start with a single first-exit time $T$ taking values in the positive integers,
possibly with mass beyond the horizon. Under a specified policy $d$, define
$S^d(j)=P(T^d>j)$ and $S^d(0)=1$. Thus $S^d(j-1)$ is the probability of an
active account at the opening of period $j$. With payment for each period
entered, the expected restricted active time is
\begin{equation}
 R^d(K)=E\{\min(T^d,K)\}
       =\sum_{j=0}^{K-1}S^d(j).
 \label{exec17:s5-restricted-time}
\end{equation}
The identity follows account by account from
$\min(T,K)=\sum_{j=0}^{K-1}1\{T>j\}$ before taking expectations. This convention
earns a period's margin even if exit occurs at its end. A different payment
clock needs its own reward function. In particular, the same numerical area
must not be priced as both opening-period activity and completed-period
activity.

Suppose the observation and identification assumptions support the policy
survival probabilities only through period $m<K$. This may be an administrative
follow-up limit, rather than random censoring that an observation model could
remove. A fitted survival function beyond $m$ then supplies an extrapolation
assumption. Increasing simulation draws does not add observations at those
times. Without a tail restriction, monotonicity of survival gives the sharp
discrete bounds
\begin{equation}
 R^d(m)+S^d(m)\ \leq\ R^d(K)\ \leq\
 R^d(m)+(K-m)S^d(m).
 \label{exec17:s5-tail-bounds}
\end{equation}
To prove the lower bound, expand the remaining terms as
$S^d(m)+S^d(m+1)+\cdots+S^d(K-1)$ and use nonnegativity. To prove the upper
bound, bound every one of these $K-m$ terms by $S^d(m)$. The known term
$S^d(m)$ cannot be discarded. The continuous-time argument that allows all
remaining activity to disappear immediately after an observed time would
miss this whole discrete period.

Both endpoints are attainable. Keep the observed distribution through $m$
unchanged. For the lower endpoint, make all accounts still active at $m$ exit
at the end of period $m+1$. For the upper endpoint, keep them active through
the entire requested horizon. Mixing these two continuation laws preserves
every observed survival probability and fills the interval. The result is an
identified set conditional on the supported earlier survival law, not a
confidence interval for a point target. Sampling uncertainty about that earlier
law is a further problem.

For an exact six-period example, take the observed policy survival sequence
at times $0,1,2,3$ to be $(1,0.8,0.65,0.5)$ and the reference sequence to be
$(1,0.9,0.8,0.6)$. Then
\begin{align}
 R^d(6)&\in[2.95,3.95],& R^0(6)&\in[3.30,4.50],\notag\\
 R^d(6)-R^0(6)&\in[-1.55,0.65].
 \label{exec17:s5-tail-example}
\end{align}
These contrast endpoints are sharp when there is no additional cross-policy
restriction: select the lower continuation for one policy and the upper for
the other. Such potential event-time marginals can coexist, for example by
constructing each from a common uniform variable and its inverse distribution
function. A scientifically justified individual monotonicity restriction
would be additional information and could narrow the set.

At USD~20 margin per period entered and USD~3 upfront policy cost, the
incremental retention-margin value is therefore in $[-34,10]$ USD per eligible
account. This example prices activity alone; it does not silently include a
default loss or a repayment penalty. Those would require the corresponding
joint event probabilities. A decision to roll out the policy cannot be
justified by selecting the favorable end of this set. A bank with a more
restrictive accepted tail model could make a different, explicitly conditional
decision.~\citep{exec17_s5_math}

There are several substantive remedies. Shortening the requested horizon to
the supported period changes the business question but removes this tail
ambiguity. Following a public or generated cohort longer can demonstrate how
a proposed analysis would use additional follow-up, without supplying the
bank's missing future. An externally justified survival envelope or a specified
parametric tail gives conditional bounds or a point extrapolation. The
report should show how the monetary conclusion changes across those choices.
An automatic default horizon chosen from the largest event in the current
sample answers a fourth question, whose end point varies with the sample.
That can be useful internally, but it does not establish comparability between
policies, software engines, or reporting vintages.

The software interface consequently needs an explicit requested horizon,
an observed-support horizon, an event convention and a monetary payment clock.
A response can legitimately contain a point estimate on the supported range
and a set beyond it. It must distinguish extrapolated uncertainty conditional
on a tail model from uncertainty that includes the choice of that model.
Returning a finite number at every requested time does not settle the
identification question.

\section{Why a late hazard can reverse a policy ranking}
\label{exec17:s5-selection}

Consider two equally weighted baseline risk strata, high and low, and no
competing exit or censoring. Within each stratum an active account defaults
with a constant probability per period. Under the reference policy these
probabilities are $0.8$ and $0.1$; under the intervention they are $0.4$ and
$0.05$. One can generate the potential histories using independent period
uniforms within each stratum. Randomized policy assignment, consistency and
the stated event clock then suffice for this illustrative policy comparison.
Its reversal does not require an unmeasured common cause of action and default.

Write $h_b^d$ for the hazard in baseline stratum $b$ and $\pi_b=1/2$ for its
baseline proportion. Marginal survival and the opening-period risk composition
are
\begin{align}
 S^d(j)&=\sum_b\pi_b(1-h_b^d)^j,\notag\\
 w_b^d(j-1)&=\frac{\pi_b(1-h_b^d)^{j-1}}{S^d(j-1)},\notag\\
 h^d(j)&=\sum_bw_b^d(j-1)h_b^d
        =1-\frac{S^d(j)}{S^d(j-1)}.
 \label{exec17:s5-mixture-hazard}
\end{align}
The two marginal hazards average over different survivor populations. By
period three, the reference high-risk share is $4/85$, whereas the policy
high-risk share is $144/505$. Protecting the high-risk accounts keeps more of
them in the later policy risk set. Comparing those risk sets as if they were
the same baseline population confuses composition with a policy's cumulative
consequence.

The exact results make the reversal visible:
\begin{center}
\begin{tabular}{lrr}
\hline
Quantity & Reference & Policy\\
\hline
$S(1)$ & $0.55$ & $0.775$\\
$S(2)$ & $0.425$ & $0.63125$\\
$S(3)$ & $0.3685$ & $0.5366875$\\
Period-three marginal hazard & $113/850$ & $1513/10100$\\
Three-period default risk & $0.6315$ & $0.4633125$\\
Restricted active periods & $1.975$ & $2.40625$\\
\hline
\end{tabular}
\end{center}
The period-three hazards are approximately $0.13294$ and $0.14980$. Thus the
policy has the larger late marginal hazard, although every stratum has a
smaller hazard under that policy at every time and marginal survival is larger
throughout. The marginal hazard ratio is a valid description of its two
selected risk sets. It is insufficient as a stand-alone ranking of the
baseline policy decision.

For a portfolio entering together, suppose each period entered earns USD~20,
default by period three loses USD~1,000, and the policy costs USD~3 upfront.
The incremental expected value is
\begin{equation}
 20(2.40625-1.975)
 -1000(0.4633125-0.6315)-3
 =173.8125\ \mathrm{USD}.
 \label{exec17:s5-selection-value}
\end{equation}
This ranking uses the same initial population and the cumulative event
probabilities appropriate to its financial objective. It remains conditional
on the stipulated margins and losses. Changing those amounts is a sensitivity
analysis of the objective, while changing how histories are selected is a
change to the population or causal model.~\citep{exec17_s5_math}

This example motivates a practical output hierarchy. Preserve stratum-specific
and marginal hazards as diagnostics, but also report the policy survival or
cumulative incidence on a common baseline population, the restricted active
time and their joint uncertainty. If the bank wants an effect among accounts
that would survive under both policies, that is a principal-stratum question
requiring additional assumptions about joint potential histories. Neither
conditioning each fitted curve on its own observed survivors nor adjusting
the plot legend identifies that population.

The same logic affects model assessment. A regression can accurately predict
late defaults among observed reference survivors and still be a poor model
for the histories retained under another policy. Checking calibration on the
observed late risk set is informative about that set. Transporting it to a
policy-induced risk set requires overlap and stability of the relevant
conditional transition law. Flexible prediction and causal transport should
therefore appear as separate arguments in the proposal.

\section{Competing exits and the order of a simulator}
\label{exec17:s5-ordered}

Now distinguish default $D$ and early repayment $P$ as mutually exclusive
first events within each period. Among accounts event-free at the opening,
let $\lambda_{D,j}$ and $\lambda_{P,j}$ be their probabilities conditional on
the permitted history and intervention. They obey
$\lambda_{D,j}+\lambda_{P,j}\leq1$. With constant probabilities in a homogeneous
example,
\begin{align}
 S(j)&=(1-\lambda_D-\lambda_P)^j,\notag\\
 F_D(K)&=\sum_{j=1}^{K}S(j-1)\lambda_D,&
 F_P(K)&=\sum_{j=1}^{K}S(j-1)\lambda_P,\notag\\
 S(K)+F_D(K)+F_P(K)&=1.
 \label{exec17:s5-first-events}
\end{align}
Here $S$ means free of both first events. In contrast,
$1-F_D=S+F_P$ includes accounts that have already repaid. It cannot be
integrated as if those repaid accounts were still earning the active-account
margin. Treating prepayment as censoring would target another quantity and
needs a separate scientific interpretation; it does not automatically
describe eliminating prepayment while keeping every other mechanism fixed.

Some simulators implement within-period draws sequentially. If prepayment is
drawn first with probability $p_P$ and default is then drawn among those
without that prepayment with probability $p_D$, the induced unconditional
first-event probabilities are
\begin{equation}
 \lambda_P=p_P,\qquad \lambda_D=(1-p_P)p_D,
 \qquad P(\text{neither})=(1-p_P)(1-p_D).
 \label{exec17:s5-ordered-conversion}
\end{equation}
Consequently an adapter from mutually exclusive first-event probabilities
must set $p_P=\lambda_P$ and
$p_D=\lambda_D/(1-\lambda_P)$ when $\lambda_P<1$. The constraint on the
two first-event probabilities ensures $p_D\leq1$. If $\lambda_P=1$, feasibility
requires $\lambda_D=0$ and the unused conditional default probability does not
determine any outcome. Reversing the order requires the converse
factorization. Multiplying two arbitrary competing hazards without defining
their conditioning populations supplies neither factorization.

The conversion is elementary but consequential. Consider two periods. Set
$(\lambda_D,\lambda_P)=(0.1,0.1)$ for reference and $(0.12,0.6)$ for policy.
Although the policy's default probability among initially event-free accounts
is larger each period, more of its accounts prepay before a later default is
possible. The resulting $(F_D,F_P,S)$ vectors at period two are respectively
$(0.18,0.18,0.64)$ and $(0.1536,0.768,0.0784)$. The policy has the smaller
default cumulative incidence and the shorter active relationship.

Price each entered period at USD~20, default at a USD~1,000 loss, prepayment
at a USD~50 loss, and the intervention at USD~3 per initial account. The
joint objective is
\begin{equation}
 V^d(K)=20R^d(K)-1000F_D^d(K)-50F_P^d(K)-c_d.
 \label{exec17:s5-event-value}
\end{equation}
Using $c_0=0$ and $c_d=3$ gives reference value $-153$ USD and policy value
$-169.4$ USD. The incremental value is $-16.4$ USD despite a smaller default
incidence. Default avoidance alone does not establish a desirable servicing
policy when relationship duration and competing exits carry monetary value.

For the ordered simulator the correct default probabilities after a
non-prepayment draw are $1/9$ for reference and $0.3$ for policy. Suppose an
adapter instead passes the policy's unconditional $0.12$ directly as $p_D$,
while setting $p_P=0.6$. The induced period default probability is then
$0.4\times0.12=0.048$. At period two the wrong policy calculation returns
\begin{align}
 F_D^{\rm wrong}&=0.064896,&F_P^{\rm wrong}&=0.8112,\notag\\
 R^{\rm wrong}&=1.352,&V^{\rm wrong}&=-81.416\ \mathrm{USD}.
 \label{exec17:s5-ordered-error}
\end{align}
Compared with the correctly translated reference, the apparent gain becomes
$71.584$ USD. An omitted conditional-probability conversion has reversed the
decision. The example does not identify a defect in an upstream package: the
error is the deliberately incompatible meaning supplied at its boundary.
\citep{exec17_s5_math}

The inspected simulation path of pygformula draws a competing event before
the target event and marks the target event missing after that competing
exit. Its aggregation multiplies the target-event probability by the
no-competing-event probability and the preceding survival products. The
stored object named survival is one minus accumulated target incidence;
it is therefore distinct from event-free survival when competing exits are
present. These are claims about the inspected original simulation and
aggregation paths, not about newly fitted models or all package releases.
The model-fitting risk sets were not newly inspected for this comparison,
so the simulation order alone cannot establish how every fitted probability
is conditioned.~\citep{v13_sw_pygformula,exec17_s5_paths}

An integration should first check the complete tuple: opening risk set,
within-period order, event codes, probability conditioning, transition to an
absorbing state, and the meaning of every accumulated output. A generated
two-period law is a useful exact semantic test before statistical fitting.
It can verify that a proposed wrapper conserves probability and obtains
Equation~\eqref{exec17:s5-event-value} under declared inputs. It does not
establish that the observational data identify those inputs, or that a
fitted observational model is calibrated under intervention.

\section{Modifying a natural action and redrawing an action}
\label{exec17:s5-natural}

A dynamic servicing policy may modify the action that would naturally be
chosen: shorten a delay, reduce a proposed payment, or add support when the
proposed contact intensity lies in a specified range. This is more specific
than drawing a new action from a distribution with the same conditional
marginal frequencies. At later times the natural action is evaluated under
the history produced by earlier interventions. Unobserved features of the
action mechanism can connect that later natural action to an earlier one,
even when there is no unmeasured confounding of the ordinary joint treatment
regime and outcome.

The longitudinal modified treatment policy analysis distinguishes a natural
action policy from its stochastic-intervention counterpart. Under its
structural model, standard sequential randomization separates a current
treatment innovation from later covariate and outcome innovations. The
stronger condition also separates it from future treatment innovations.
Theorem~1 identifies the redrawn stochastic intervention under the standard
condition and the natural modified policy under the stronger condition,
together with support. The two targets must therefore be named before
reusing the same observed-data recursion.~\citep{lmtp2023}

An exact two-time construction shows why this distinction is necessary.
Let $U$ and $V$ be independent fair bits, with no measured covariates. In both
worlds $A_1=U$ and the structural outcome is $Y=A_2$. The second natural action
is generated by
\begin{equation}
 \begin{array}{ll}
 \mathcal W_0:& A_2=V,\\
 \mathcal W_1:&
 A_2=\begin{cases}V,&A_1=U,\\1,&A_1\ne U.\end{cases}
 \end{array}
 \label{exec17:s5-two-worlds}
\end{equation}
In ordinary observations $A_1=U$, so the second branch of $\mathcal W_1$ is
never used. In each world all four $(A_1,A_2)$ cells have probability $1/4$,
and $Y=A_2$. Thus the complete observable law is identical, with positive
support for both actions at each observed history. This is equality of the
whole joint distribution, not merely equality of observed outcome means.

Apply the policy $d_1(a)=1-a$, followed by $d_2(a)=a$: flip the first natural
action and leave the second natural action as generated under the changed
history. In $\mathcal W_0$, the second action remains $V$ and
$E(Y^{\bar d})=1/2$. In $\mathcal W_1$, the first action becomes $1-U$;
the second action mechanism then takes its other branch and sets $A_2=1$.
Its modified-policy mean is $1$. The natural-action effect is not identified
by the common observed law under the weaker restrictions of this construction.

For the redrawn intervention, the observed conditional action law at both
times is a fair bit. Pushing those conditional laws through the flip and
identity maps still gives fair bits; drawing from them under the evolving
history yields outcome mean $1/2$ in both worlds. This is a different
intervention. It is identified by the common observed law under the stated
standard condition. Calling it a numerical approximation to the natural
policy would obscure a difference in counterfactual meaning.

The construction satisfies the standard innovation condition used above:
covariate and outcome structural innovations can be constants, and the
outcome under a joint assigned sequence is simply the assigned second
action. The second treatment innovation in $\mathcal W_1$, however, contains
the earlier $U$. Strong independence between the first and later treatment
innovations fails. In $\mathcal W_0$ independent $U,V$ suffice. Support also
holds in the example; the failure does not arise from a missing observed
action cell. These details explain why the witness illustrates the theorem's
assumption distinction rather than contradicting its identification result.
\citep{lmtp2023,exec17_s5_math}

In banking terms, a latent servicing preference that influences several
successive proposed actions can matter to a policy that modifies those
proposals. Observing all outcome confounders for ordinary assigned regimes
does not automatically reveal that proposal mechanism. This interpretation
is a possible mechanism, not an assertion about a bank's data. A proposal
should ask whether the natural recommendation can be recorded, whether its
history-dependent generation can be justified, and whether a randomized
redraw is operationally acceptable. If the business instead intends a fresh
randomized allocation at every time, the stochastic target may be exactly
the intended one.

There are consequently several defensible responses to uncertain natural
action dependence: record and experimentally vary the recommendation
mechanism; impose and defend the stronger structural restriction; redesign
the intervention as a redrawn policy; or retain an explicitly incomplete
identification claim. An estimator's robustness properties concern an
identified observed-data functional under its assumptions. They cannot
choose among observationally identical structural worlds with different
natural-policy outcomes.

\section{A supported shift, a cap, and the boundary of the density ratio}
\label{exec17:s5-pushforward}

Continuous treatment policies create a second boundary that is easy to hide in
an API. Let $A\sim\operatorname{Uniform}[0,1]$ and consider the fixed map
\begin{equation}
 d(a)=\begin{cases}a+1/4,&0\leq a\leq3/4,\\a,&3/4<a\leq1.\end{cases}
 \label{exec17:s5-shift}
\end{equation}
The map is not globally one-to-one. Its image has two inverse branches on
$[3/4,1]$: the identity branch and the shifted branch. The pushforward density
relative to Lebesgue measure is therefore
\begin{equation}
 g_d(y)=\begin{cases}
 0,&0\leq y<1/4,\\
 1,&1/4\leq y<3/4,\\
 2,&3/4\leq y\leq1.
 \end{cases}
 \label{exec17:s5-push-density}
\end{equation}
The value two is not a probability exceeding one. It is the sum of the masses
from two preimages. Dropping one inverse branch changes the intervention law
and its likelihood ratio. Integrating the density gives one, and direct
integration gives
\begin{equation}
 E\{d(A)\}=11/16,\qquad E\{d(A)^2\}=25/48,\qquad
 \operatorname{Var}\{d(A)\}=37/768.
 \label{exec17:s5-push-moments}
\end{equation}

Now compare the cap $d_c(a)=\min(a+1/4,1)$. It maps the
interval $(3/4,1]$ to the single point one, creating an atom of mass $1/4$.
The observed continuous treatment law gives that point probability zero, so
the intervention law is not absolutely continuous with respect to the observed
law. A density-ratio estimator defined only with respect to the observed
Lebesgue density cannot represent this target by an ordinary Radon--Nikodym
ratio. The endpoint belongs to the topological support $[0,1]$; the problem
is the stronger absolute-continuity requirement. A large finite ratio and a
singular atom are different mathematical objects.

The distinction matters even if an outcome regression is available. Suppose
$m(a)=E(Y\mid A=a)$ is known only as a version up to Lebesgue-null sets and
the outcome is bounded in $[0,1]$. The value at the cap point is then
unrestricted by the observed regression. Suppose consistency and
exchangeability identify the causal response mean almost everywhere on that
continuum, without imposing endpoint continuity. The target mean then lies in
\begin{equation}
 \left[\int_0^{3/4}m(a+1/4)\,da,
       \int_0^{3/4}m(a+1/4)\,da+1/4\right]
 = [15/32,23/32]
 \label{exec17:s5-cap-set}
 \end{equation}
for the exact authored regression $m(a)=a$ almost everywhere on the observed
continuum. The interval is sharp. Let $E$ be uniform on $[0,1]$, independent
of $A$, and define $Y^a=1\{E\leq m_u(a)\}$, where $m_u(a)=a$ for $a<1$ and
$m_u(1)=u$ with $u\in[0,1]$. Every such world has the same complete
observed $(A,Y)$ law, because $P(A=1)=0$. Its capped mean is
$15/32+u/4$. The two endpoint worlds and their intermediate values therefore
attain the entire set while preserving exchangeability and consistency.

If the outcome is worth USD~100 at its unit value, the baseline is USD~50 and the
intervention costs USD~3, the net value set is
$[-49/8,151/8]$ USD. A continuous extension of $m$ to the endpoint would
force $u=1$ in this example. A structural model for the atom supplies a
different explicit restriction. Redesigning the intervention to avoid the
atom changes the policy instead. These remedies address distinct scientific
questions; they must not be interchanged by an automatic clipping rule.
\citep{lmtp2023,exec17_s5_math}

Identification under continuity is also distinct from regular estimation.
A value defined at a single continuous treatment point can depend on local
behavior that the nonparametric observed-data model does not estimate at
root-$n$ rate. The LMTP paper's threshold example describes this
non-pathwise-differentiability and restricts its efficiency theorem to
piecewise smoothly invertible maps. Continuity can resolve the choice of an
endpoint version without proving the square-integrable influence function
required by that theorem. A parametric model, quantitative smoothness
restriction, mixed treatment law with a genuine endpoint atom, or a smoothed
policy would need its own estimation argument. The cap is operationally
meaningful; the unrestricted continuous-density theorem is simply not its
automatic uncertainty guarantee.~\citep{lmtp2023}

\section{The derivative of a fixed modified-action map}
\label{exec17:s5-derivative}

For the supported fixed map in Equation~\eqref{exec17:s5-shift}, write
$g(a\mid W)$ for the observed conditional treatment density,
$r(a,W)=g_d(a\mid W)/g(a\mid W)$, and $m(a,W)$ for the conditional outcome
regression. Here $W$ denotes baseline information available before the action.
Assume the relevant change of variables is valid, the map is fixed, and the
score below has finite second moment. In a one-time setting, consistency,
exchangeability given $W$, and the required absolute continuity identify the
following causal mean from the observed law. The one-time target is
\begin{equation}
 \theta(P)=E_W\int m(d(a,W),W)g(a\mid W)\,da.
 \label{exec17:s5-lmtp-target}
\end{equation}
One can derive its gradient directly. For a smooth dominated submodel write
the score as $s=s_W+s_A+s_Y$, with
$s_W=E(s\mid W)$,
$s_A=E(s\mid A,W)-E(s\mid W)$, and
$s_Y=s-E(s\mid A,W)$. The three scores have mean zero in their respective
conditioning populations. Differentiation of the outer baseline law, the
assignment law inside the integral, and the conditional outcome regression
gives
\begin{align}
 \dot\theta={}&E\{(M(W)-\theta)s_W\}\notag\\
 &+E\{(m(d(A,W),W)-M(W))s_A\}\notag\\
 &+E_W\int g_d(a\mid W)
       E\{(Y-m(a,W))s_Y\mid A=a,W\}\,da ,
 \label{exec17:s5-path-parts}
\end{align}
where $M(W)=E\{m(d(A,W),W)\mid W\}$. The last line uses the pushforward
change of variables, including every inverse branch, to move the derivative
of $m(d(a,W),W)$ onto the observed treatment scale. Replacing $g_d$ by $rg$
makes that line the expectation of $r(A,W)(Y-m(A,W))s_Y$.
The fixed-map path derivative thus has the mean-zero influence function
\begin{equation}
 D(P)(O)=r(A,W)\{Y-m(A,W)\}+m(d(A,W),W)-\theta(P).
 \label{exec17:s5-lmtp-if}
\end{equation}
Each summand recovers its corresponding score covariance in
Equation~\eqref{exec17:s5-path-parts}; all cross-covariances vanish by
conditional centering. This establishes the derivative in the stated
observational model. The primary LMTP theorem supplies the longitudinal
generalization under its additional conditions.~\citep{lmtp2023}

The map is fixed independently of $P$, although its induced density $g_d$
changes when $g$ changes. To isolate the distinction, define a second
functional using an external intervention density $q_0(a\mid W)$ held fixed
along the submodel. At the reference law take $q_0=g_d$, so the two targets
have the same numerical value at that law. Under perturbation they need not.
For the external-$q_0$ functional the middle line of
Equation~\eqref{exec17:s5-path-parts} is absent and
\begin{align}
 M(W)&=\int m(a,W)q_0(a\mid W)\,da,\notag\\
 D_{\mathrm{frozen-q}}&=r(A,W)\{Y-m(A,W)\}+M(W)-\theta.
 \label{exec17:s5-frozen-if}
\end{align}
The observed assignment law may still vary; it is the intervention law that
is externally frozen in this comparison. A known randomized observational
assignment defines a restricted statistical model and a different
projection argument. Neither interpretation should be chosen merely from
the fact that a software routine accepts supplied probabilities.

The extra random term in the observational fixed-map score is
$m(d(A,W),W)-M(W)$. Its covariance with the outcome-residual component
vanishes after conditioning on $(A,W)$, and its covariance with the
baseline component vanishes after conditioning on $W$. Consequently
\begin{equation}
 \operatorname{Var}(D)-\operatorname{Var}(D_{\mathrm{frozen-q}})
 =E\!\left[\operatorname{Var}\{m(d(A,W),W)\mid W\}\right]\geq0.
 \label{exec17:s5-assignment-variance}
\end{equation}
This is a comparison of the two specified influence functions at the reference
law, not a claim that their estimators answer the same question away from it.

The exact check uses the shift in Equation~\eqref{exec17:s5-shift} and
$m(a)=a$, with no baseline variation. Along the normalized assignment path
$g_\epsilon(a)=1+\epsilon(2a-1)$ for $|\epsilon|<1$, direct integration gives
\begin{align}
 \theta(\epsilon)
 &=\int_0^{3/4}(a+1/4)\{1+\epsilon(2a-1)\}\,da\notag\\
 &\quad+\int_{3/4}^1 a\{1+\epsilon(2a-1)\}\,da
   =\frac{11}{16}+\frac{23}{192}\epsilon .
 \label{exec17:s5-assignment-path}
\end{align}
Thus $\dot\theta(0)=23/192$ and $\theta(1/4)=551/768$. The covariance
$E[\{d(A)-11/16\}(2A-1)]$ reproduces the derivative; an outcome-residual
score alone gives zero and misses it. Under the external frozen intervention
law, with $m$ held fixed, zero is instead the correct derivative. The check
separates the scientific targets before comparing estimated standard errors.

For a Bernoulli outcome with mean $A$, the residual score variance is
\[
 \int_{1/4}^{3/4}a(1-a)\,da+
 4\int_{3/4}^1a(1-a)\,da=\frac7{32}.
\]
The factor four is the square of the density ratio two, not a second
contribution to the probability mass. Adding the assignment variance
$37/768$ from Equation~\eqref{exec17:s5-push-moments} gives full score
variance $205/768$. For $n$ independent observations an asymptotically
linear estimator with this score has leading variance $205/(768n)$,
subject to its nuisance-rate and remainder conditions. The arithmetic
checks those components but does not establish those conditions for a
fitted learner or an actual bank. The monetary outcome scale multiplies
the score linearly and its variance quadratically.~\citep{exec17_s5_math}

\section{Translating software semantics before comparing estimates}
\label{exec17:s5-software}

The common package names in a longitudinal workflow conceal different data
contracts. The inspected ltmle documentation uses wide, time-ordered nodes;
an absorbing event is encoded as event one, censoring is a separate factor,
and later nodes are cleaned according to that event and censoring history.
The detailed source analysis in Section~\ref{exec23:ltmle-regimes} establishes
an additional admission requirement: absorbing outcomes are validated before
post-event records are cleaned. Its documented bounds act on cumulative
treatment probabilities. A continuous outcome range can also be
truncated or scaled. These choices determine the observed target before a
targeting step is run.~\citep{sw_ltmle}

The lmtp documentation uses chronological treatment, covariate and outcome
nodes. A supplied \texttt{shifted} data frame overrides the \texttt{shift}
function; it is not an additional rule applied after that function. The
finite-memory parameter specifies how many previous time points are used
for estimation. That documented modeling choice does not establish that a
bank's history can be truncated without changing its conditional law.
Cluster identifiers, folds and sampling weights affect the fitted calculation.
Its returned density ratios are noncumulative. Multiplying them to form the
cumulative ratios required by a stated score is valid. Treating each
noncumulative ratio as an already cumulative quantity, or multiplying an
already cumulative output again, changes that score.~\citep{sw_lmtp}

For a two-step exact probability law with incremental denominators
$(1/2,2/5)$, the incremental ratios are $(2,5/2)$ and their cumulative
products are $(2,5)$. If a floor of $1/4$ is applied to the cumulative
denominator, the final weight is $4$; if the floor is applied to each
incremental denominator before multiplying, the final weight is $5$. Both
procedures avoid an infinite value, but they do not implement the same
numerical calculation. The floor level and location belong in the estimator
specification. A fixed floor can leave bias for the original target, while
a vanishing floor needs rate and moment arguments. Clipping a denominator
does not automatically define a different normalized intervention law.
If a deliberately changed target is intended, that target needs its own
definition and identification argument.~\citep{exec17_s5_math}

The stremr manual presents a long-data interface with subject and time fields
and separates treatment, monitoring and censoring models. Its `fitTMLE` path
wraps sequential g-computation with targeting. The survtmle manual fixes a
binary treatment, integer event times and an event code for censoring; its
default horizon is derived from the observed event time. It documents mean and
hazard methods, covariance and subject-level influence vectors, and warns that
finite-sample convergence is not guaranteed. The prose and argument material
also use different descriptions of a tolerance scaling: the displayed call
uses $1/\sqrt n$, while the descriptive paragraph names $1/n$.
An adapter must inspect the selected executable path and record the actual
argument, rather than resolve that discrepancy by choosing whichever
default sounds preferable. Reaching the maximum iteration count and meeting
the empirical-score tolerance are different termination events.
Section~\ref{sec:exec23-longr-workers} follows the selected stopping loops,
optimizer fallbacks and returned influence contributions. Its findings make
final residuals and effective update counts part of the qualification, without
establishing native coverage or silently settling every tolerance convention.
\citep{v12_sw_stremr,v12_sw_survtmle}

The survtmle horizon illustrates why defaults need scientific review.
The saved API uses the maximum failure time among noncensored event records,
not the end of all recorded follow-up. Two samples with identical requested
six-period objectives can therefore have different implicit output times.
Passing an explicit common horizon removes that software difference, but
does not create positivity of observation or identify an unsupported tail.
Likewise, its covariance and subject influence vectors are useful for
post-hoc contrasts only if the requested treatment, event, horizon and
subject indexing remain attached to their entries. The interface itself
does not establish successful targeting or valid nuisance rates.

The gfoRmula API exposes risk and end-of-follow-up outcomes, censoring and
competing-event options, and custom intervention and history roles. Its names
do not by themselves establish whether an output is event-free survival,
cumulative incidence, or a different complement. The pygformula simulation
path draws competing events before target events, while its bootstrap resamples
whole histories for unique subjects, refits models, and uses common seeds for
intervention simulations. Those are useful building blocks for a prospective
comparison. They do not prove that every model-fitting risk set or every
release has the same semantics.~\citep{sw_gformula,v13_sw_pygformula,v131_sw_pygformula}

These packages also answer different design questions. A fixed binary
baseline-treatment survival analysis can be compared with survtmle's
declared input. Repeated binary decisions with censoring and an ordered
wide history fit the selected ltmle contract; a long history with a distinct
monitoring process motivates examining stremr's corresponding treatment,
monitoring and censoring components. Modified continuous action laws motivate
the lmtp map, support and ratio checks developed above. Parametric
g-computation motivates gfoRmula or pygformula when the goal includes
simulating full intervention histories. An adapter can make their results
comparable only on an intersection of targets and observation assumptions.
It cannot turn a fixed baseline treatment into an evolving decision merely
by repeating its value down a long table.

The distinction between ordinary and sequential double robustness also
belongs in this comparison. The inspected sequential-robustness paper
distinguishes consistency when either every outcome regression is correct
or every treatment mechanism is consistently estimated from consistency
when, at each time, the appropriate outcome or treatment condition holds.
Its outcome conditions concern functional specification of the recursively
defined regressions, and its formal rate and inference conditions add
further restrictions. A product should therefore identify which algorithm,
target and recursion support a robustness statement. The name TMLE, an
ensemble learner, or an empirical score near zero cannot by itself establish
the stronger per-time property. No rederivation of the earlier sequential
remainder is needed to make this software choice explicit.
\citep{sequentialdr2018}

An interoperable result should therefore carry at least: subject unit and
cluster, clock and horizon, event and censoring codes, absorbing-state rule,
policy law and whether it is natural or redrawn, treatment probability scale,
ratio orientation and truncation, target type, influence-vector convention,
failed-replicate handling, software version and license. A result lacking one
of these fields can remain a useful exploratory number, but it cannot enter a
common-target comparison without an explicit refusal or translation proof.
Package declarations currently record gfoRmula 1.1.2 (GPL-3), ltmle 1.3-0
(GPL-2) and lmtp 1.5.4 (AGPL-3.0). The inspected API bytes and package
declarations are separate evidence; metadata alone does not prove that those
bytes belong to those exact releases. The selected pygformula source is
pinned to commit \texttt{e66f787fcef2} dated 22 May 2026; the selected
stremr and survtmle records identify 0.4 at \texttt{14d15651b247} and 1.1.1
at \texttt{295de8b4be18}, respectively. These are scoped historical identities,
not a claim that all current releases were surveyed. Full dependency rights
and runtime qualification require separate inspection and trials.
\citep{exec17_s5_paths}

\section{Joint event uncertainty and the monetary objective}
\label{exec17:s5-joint}

The financial objective in Equation~\eqref{exec17:s5-event-value} depends on
three components of the same account history. They should not be treated as
independent measurements. Return to its two-period policy law, with constant
first-event probabilities $d=3/25$ for default, $p=3/5$ for prepayment, and
$s=1-d-p=7/25$ for neither event. Write $B$ for the indicator of no first-period
event, $L=1+B$ for the number of periods entered, and $I_D,I_P$ for indicators
of default and prepayment by period two. The complete subject vector is
$X=(L,I_D,I_P)^\top$. Its mean is
\begin{equation}
 E(X)=\left(1+s,\ d(1+s),\ p(1+s)\right)^\top
      =\left(\frac{32}{25},\frac{96}{625},\frac{96}{125}\right)^\top.
 \label{exec17:s5-joint-mean}
\end{equation}
This follows by conditioning on the first event: an account enters period two
exactly when $B=1$, and can default by the horizon either immediately or after
that event-free first period. No independence between the entries of $X$ has
been assumed.

The covariance has a direct indicator derivation. Since default and prepayment
are mutually exclusive first-event categories, $I_DI_P=0$ account by account.
Thus, writing $f_D=d(1+s)$ and $f_P=p(1+s)$,
\begin{align}
 \operatorname{Var}(L)&=s(1-s),\notag\\
 \operatorname{Cov}(I_D,I_P)&=-f_Df_P,\notag\\
 \operatorname{Cov}(L,I_D)
 &=E(BI_D)-s f_D=sd-sd(1+s)=-ds^2,\notag\\
 \operatorname{Cov}(L,I_P)&=-ps^2.
 \label{exec17:s5-joint-identities}
\end{align}
The remaining diagonal entries are $f_D(1-f_D)$ and $f_P(1-f_P)$.
Whole-history enumeration and these separate identities give the same matrix:
\begin{equation}
 \Sigma=\operatorname{Cov}(X)=
 \begin{pmatrix}
 126/625&-147/15625&-147/3125\\
 -147/15625&50784/390625&-9216/78125\\
 -147/3125&-9216/78125&2784/15625
 \end{pmatrix}.
 \label{exec17:s5-joint-matrix}
\end{equation}
Negative event-indicator covariance expresses exclusion of the alternative
exit. Negative duration--event covariance expresses earlier termination of
the active relationship. Both arise from the same joint history law.

For monetary coefficients $c=(20,-1000,-50)^\top$, the individual net reward
is $Z=c^\top X-3$. The cost is constant, so it affects the mean but not the
variance. The exact moments are
\begin{align}
 E(Z)&=-169.4\ \mathrm{USD},\notag\\
 \operatorname{Var}(Z)&=c^\top\Sigma c
                 =119207.04\ \mathrm{USD}^2,\notag\\
 \sum_{k=1}^3c_k^2\Sigma_{kk}&=130533.12\ \mathrm{USD}^2.
 \label{exec17:s5-joint-money}
\end{align}
Dropping the off-diagonal terms overstates the variance by
$11326.08\ \mathrm{USD}^2$ in this construction. Independence is not necessarily
an anticonservative approximation; its direction depends on the objective's
signs and the joint law. Changing those signs or changing the event mechanism
can change the direction. An integration should preserve the covariance,
rather than assume that discarding it always gives a safer answer.

For $M=1000$ independent complete trajectories drawn from this fixed policy
law, the sample-average net reward has variance $119.20704\ \mathrm{USD}^2$.
This is an exact sampling or Monte Carlo statement for those trajectories.
It is not the variance of a causal estimator fitted to $1000$ observational
subjects with estimated treatment and observation mechanisms. For that
estimator, a valid vector of subject influence contributions $\phi_i$ would
instead supply the covariance appropriate to its fitted target. Under its
own asymptotic-linear and moment conditions, the monetary influence is
$c^\top\phi_i$; a constant intervention cost contributes no influence, while
an estimated cost can contribute another component. Jointly estimated
policy and reference outputs also require their cross-policy covariance.
The covariance identity remains the same, but the random vector must be the
correct one.~\citep{exec17_s5_joint_error}

This distinction guides a substantive adoption choice. When an engine exposes
only marginal standard errors, a monetary combination may be unidentified
from its reported uncertainty even when every point estimate is usable.
Exposing subject-aligned influence vectors or a justified joint bootstrap
can repair that interface. If neither is available, one may derive a
conservative bound from additional covariance restrictions or return
unavailable joint uncertainty. Summing marginal variances asserts an
independence restriction that competing first events visibly violate here.

\section{Sampling, simulation and model error are different quantities}
\label{exec17:s5-uncertainty}

Let $D_n$ denote the full observed sample of independent subject histories,
including the records used for fitting and the empirical baseline population.
Conditional on it, let $U$ contain the simulator's baseline resampling and
transition draws. For a value estimate using $M$ simulated trajectories,
the law of total variance gives
\begin{equation}
 \operatorname{Var}(\widehat V)=
 \operatorname{Var}_{D_n}\{E(\widehat V\mid D_n)\}+
 E_{D_n}\{\operatorname{Var}(\widehat V\mid D_n)\}.
 \label{exec17:s5-total-variance}
\end{equation}
The first term is variation in the fitted baseline and observed sample; the
second is simulation noise conditional on that fitted object. If the simulator
is conditionally unbiased for a fitted-model functional $V(D_n)$, its first
conditional mean is exactly $V(D_n)$. Model misspecification can leave a
bias $E\{V(D_n)\}-V_{\rm target}$ that is not removed by increasing $M$; it
can also affect the two variance terms. Squared-error risk adds the square
of that bias to the displayed variance. A
reported Monte Carlo standard error that omits the first term is not a
confidence interval for the full causal estimand.

For a paired policy contrast, use one common uniform draw per simulated unit
and set the event indicator to $1\{U<p_d\}$. With $p_0=0.2$ and $p_1=0.3$, the
single-draw difference has mean $0.1$ and variance $0.09$, because it is one
when $U\in(0.2,0.3]$ and zero otherwise. Independent draws have variance
$0.2(0.8)+0.3(0.7)=0.37$. At $M=10{,}000$, the corresponding simulation
variances are $0.000009$ and $0.000037$. Common random numbers can improve a
contrast by coupling comparable paths, but resetting a seed alone does not
guarantee that branching histories remain aligned. The coupling rule must be
declared.

If each independent account contributes two identical rows, treating rows as
iid understates the variance of a subject-level mean by a factor of two. The
bootstrap and influence calculation must resample or retain the independent
subject unit. The pygformula bootstrap's whole-history resampling is therefore
semantically different from a row bootstrap. A future package comparison must
record this unit explicitly.~\citep{v131_sw_pygformula,exec17_s5_math}

Coverage also needs a complete denominator. In an illustrative 200-request
study, 160 requests are available, 152 are covered, so availability is $0.8$,
conditional coverage is $152/160=0.95$, and request coverage is $152/200=0.76$.
Reporting only the conditional fraction would hide the failed or unavailable
requests. Pointwise intervals for several horizons or policies are not
simultaneous bands; retaining the joint subject influence vectors permits a
declared covariance calculation, while simultaneous control needs an explicit
multiplier or resampling rule. These are reporting choices with direct business
meaning when a decision is based on the worst horizon or the best policy.

\section{A bounded transition-error envelope}
\label{exec17:s5-error}

An integration can state a conditional transition-error budget while retaining
its scientific premises. Suppose the reference and approximate processes have
the same initial history law on standard Borel spaces, use the same fixed
policy and reward functions, and differ only in transition kernels
$P_t(\cdot\mid h)$ and $\widetilde P_t(\cdot\mid h)$. Include any randomized
policy draw in the common transition description. Define total variation as
$\sup_E|P_t(E\mid h)-\widetilde P_t(E\mid h)|$, and suppose it is at most
$\varepsilon_t\in[0,1]$ uniformly over histories reachable by either process.
These are restrictions on the compared laws, not conclusions from predictive
fit on the observed histories.

Couple the initial histories identically. While their full histories agree,
use a measurable maximal coupling of the next kernels; its conditional
probability of a first split is at most $\varepsilon_t$. Once histories have
split, regard the split flag as permanent, even if current states later agree.
If $a_t$ denotes the actual probability of a split by step $t$, then
$a_{t+1}\leq a_t+(1-a_t)\varepsilon_t$. The map on the right is increasing
in $a_t$. Induction therefore bounds $a_t$ by
\begin{equation}
 q_0=0,\qquad q_{t+1}=q_t+(1-q_t)\varepsilon_t,
 \qquad q_t=1-\prod_{u=0}^{t-1}(1-\varepsilon_u).
 \label{exec17:s5-coupling-recursion}
\end{equation}
The event of unequal histories is contained in the split event; this also
bounds the total-variation distance between their history laws. A simpler
union bound, useful when an additive budget is easier to report, gives
\begin{equation}
 d_0=0,\qquad d_{t+1}\leq\min\{1,d_t+\varepsilon_t\}.
 \label{exec17:s5-tv-recursion}
\end{equation}
One may choose $d_t=\min\{1,\sum_{u<t}\varepsilon_u\}$ as a conservative
bound, with $q_t\leq d_t$. For a common running reward with range at most
$B_t$ on step-$t$ histories, the absolute
expected running-reward bias is at most
\begin{equation}
 \sum_{t=0}^{K-1}B_t q_t\ \leq\ \sum_{t=0}^{K-1}B_t d_t,
 \label{exec17:s5-running-bound}
\end{equation}
and a common terminal loss with range at most $L$ adds at most $Lq_K$,
or conservatively $Ld_K$. To see this,
condition on the common history event and use equality of coupled rewards
there; on the split event the reward difference is bounded by its range.
Summing those expected absolute differences and applying the triangle
inequality proves the stated envelope. If the reward functions themselves
are approximated, their uniform discrepancies add a separate term; the
displayed transition bound does not already include them.

For $K=3$, $\varepsilon_t=0.02$, $B_t=20$ USD and terminal loss range
$L=1000$ USD, the recursion gives
$d=(0,0.02,0.04,0.06)$, whereas the tighter coupling recursion gives
$q=(0,0.02,0.0396,0.058808)$. The additive running-margin and terminal
bounds are USD~1.20 and USD~60, totaling USD~61.20. The tighter components
are USD~1.192 and USD~58.808, totaling USD~60. The difference is a proved
consequence of counting the opportunity for a first split only while the
histories remain coupled.~\citep{exec17_s5_joint_error}

The tighter envelope is attainable under these restrictions. Let the
reference stay in state zero. Let the approximation, while in state zero,
move to an absorbing state one with probability $\varepsilon_t$; otherwise
it stays zero. At state one define both kernels to stay there. The two
kernels differ in total variation by exactly $\varepsilon_t$ at state zero
and agree at state one. The approximate probability of state one at time
$t$ is $q_t$, whereas the reference probability is zero. Choose the common
running reward $B_t1\{\text{state}=1\}$ and terminal reward
$L1\{\text{state}=1\}$. Their expected difference is exactly
$\sum_tB_tq_t+Lq_K$, attaining all components together. This abstract witness
establishes sharpness of the envelope over the declared kernel/reward class.
It does not imply that the stipulated budgets hold for a fitted bank model.
\citep{exec17_s5_joint_error}

A bank can replace these
uniform budgets with validated conservative budgets only after obtaining
appropriate transition audits. A small average error on frequent histories
does not imply this uniform bound on rare but financially important histories.
Likewise, a terminal value range should include all competing event outcomes,
not only defaults, if the policy objective prices them jointly.

The envelope offers a useful admission test for an implementation. The
adapter must expose the horizon, transition decomposition, reward range,
stated TV or coupling budget, and whether the bound includes model fitting,
simulation, and observation error. More precisely, the displayed envelope
bounds the expected reward under two fixed laws; finite Monte Carlo error
and fitted-law sampling uncertainty need the separate arguments above.
For a policy-reference contrast, bounding each policy's law separately and
adding the two envelopes is valid; a smaller paired envelope requires a
justified common coupling. If any budget is unavailable, the system
can return a conditional point estimate with a declared numerical component,
but it should not label the resulting interval a full causal uncertainty
interval. Refusal is a valid output when the requested guarantee exceeds the
declared budget.

\section{A common-target public prototype and adoption routes}
\label{exec17:s5-adoption}

The practical next step is a public or generated comparison with a common
target, not a claim that six packages are interchangeable. Freeze a finite
baseline cohort law, an integer event clock, explicit default/repayment/censor
codes, two policies, a supported horizon, monetary rewards, and a complete
seed and environment manifest. Admit an engine only after exact two-period and
single-shift references reproduce the probability, pushforward and reward
checks above. A candidate that cannot express an atom, competing first event,
natural-versus-redrawn choice, or subject-level resampling is not silently
coerced into the target; it receives a typed incompatibility record.

For each admitted engine, register the input schema, version and license,
training folds, policy specification, simulation coupling, number of baseline
and Monte Carlo draws, interval or band method, failed-replicate rule and
runtime. Compare the policy contrast, restricted active time, each first-event
incidence, joint covariance, availability, width, numerical error and the
resulting monetary decision. Keep model bias as a separate unknown. A generated
potential-outcome law establishes the truth of that construction only. Public
data can exercise the observation and clock pipeline, but cannot establish a
bank's unmeasured assignment truth or accepted loss threshold.

Adoption then has four defensible routes. Direct adoption is suitable when an
engine's target, event semantics, license and runtime are accepted. A thin
adapter is preferable when the estimator is sound but its schema or result
metadata differ. Porting is justified only when the missing capability is
material, the license permits it, and a reference implementation can be
reproduced. Rejection is appropriate for an unsupported target, unresolvable
license or dependency risk, missing uncertainty meaning, or an output that
cannot preserve the bank's information clock. These are proposal decisions,
not claims that the underlying packages are defective.

The shared library should expose a typed analysis manifest and a refusal
object alongside estimates. The CLI can validate the manifest, run exact
semantic fixtures, and emit a versioned result bundle. An MCP surface can
request validation, explanation, and comparison while preserving the same
manifest and audit record; it should not infer a causal diagram or silently
choose a horizon from an unqualified natural-language request. LLM or agent
assistance can propose mappings and diagnostics, but a human or registered
rule must approve the target, assumptions, source clock and policy. The
longitudinal engine should return an identified set or an incompatibility when
the tail, atom, event order or natural-policy condition is unsupported.

A prospective public/generated pilot can test these translations before an
expensive bank experiment. Its adoption decision should distinguish an exact
semantic match from statistical qualification, complete dependency rights,
accepted operating behavior and economic usefulness. Authentic assignment and
observation records, feasible policy versions, independent review and accepted
bank loss thresholds remain necessary for those later decisions.
\citep{exec17_s5_design}



% SOURCE docs/release/2026-10-29/source/docs/writing-repair/2026-10-29/src/16b19f74b2e54c80.tex CHAR 0
\section{Transition error along the histories that a policy actually visits}
\label{exec17:s5-occupied}

The uniform envelope above has a useful virtue: it remains valid without
knowing where a transition error occurs. Its price is that a large loss range
can be charged at every possible first split. A servicing comparison can be
more informative when it preserves both the histories visited under the
policy and the financial consequences of the particular transitions that
differ. This requires additional information. It cannot be obtained by
renaming an average predictive error as a uniform transition budget.

Consider a finite horizon of $K$ transitions and a finite state $X_t$ containing
the information needed by a fixed policy. The full finite history may be used
as the state; compressing it into a balance and a delinquency band requires a
separate sufficiency assumption. Let $P_t$ be the reference transition kernel
and $Q_t$ its approximation after applying the same policy. They start from
the same distribution $\mu_0$ and use the same running rewards $r_t(x)$ and
terminal reward $g(x)$. Thus the comparison concerns model error for one
policy. It is not itself the causal effect of switching policies. Put
\begin{equation}
 \begin{split}
 v^Q_K(x)&=g(x),\\
 v^Q_t(x)&=r_t(x)+\sum_y Q_t(y\mid x)v^Q_{t+1}(y),\\
 V(Q)&=\sum_x\mu_0(x)v^Q_0(x),\\
 \mu^P_t(x)&=\Pr_P(X_t=x).
 \end{split}
 \label{exec17:s5-occupied-values}
\end{equation}
The corresponding $P$ recursion defines $v^P$ and $V(P)$. The rewards can
include active-account margin and terminal indicators for mutually exclusive
default and prepayment. The latter are sufficient for undiscounted event
losses. Discounted payments or event-time costs need the corresponding
time/state reward representation; an unchanged terminal indicator cannot
silently account for them.

Define the signed continuation residual
\begin{equation}
 b_t(x)=\sum_y\{P_t(y\mid x)-Q_t(y\mid x)\}v^Q_{t+1}(y).
 \label{exec17:s5-occupied-residual}
\end{equation}
Subtracting the two backward recursions gives
$v^P_t-v^Q_t=P_t(v^P_{t+1}-v^Q_{t+1})+b_t$. Average against $\mu^P_t$.
The first term becomes the average of the next value difference against
$\mu^P_{t+1}$. Repeating this step cancels every intermediate difference;
the terminal difference is zero because $g$ is common. Hence the exact identity
is
\begin{equation}
 V(P)-V(Q)=\sum_{t=0}^{K-1}\sum_x\mu^P_t(x)b_t(x).
 \label{exec17:s5-occupied-identity}
\end{equation}
No small-error expansion is used. The approximation's continuation value is
weighted by occupancy under the reference process. Replacing that occupancy
by the approximation's occupancy in this formula is generally incorrect.
There is a second exact identity with $Q$ occupancy and $P$ continuation
values, obtained by reversing the subtraction. It does not make the unknown
reference continuation values available from an approximate simulator.

For a row $(t,x)$, let $\epsilon_t(x)$ bound its total variation and let
$R_t(x)$ be the range of $v^Q_{t+1}$ over a set containing the support of the
signed row difference. That signed difference has total mass zero. Its
positive and negative masses are equal to its total variation, so subtracting
the smallest continuation value proves
\begin{equation}
 |b_t(x)|\leq \epsilon_t(x)R_t(x),\qquad
 |V(P)-V(Q)|\leq
 \sum_{t,x}\mu^P_t(x)\epsilon_t(x)R_t(x).
 \label{exec17:s5-occupied-envelope}
\end{equation}
Using the full next-state space gives a valid range even when the signed
support is unknown. A smaller range is valid only when the excluded
transitions are known to have zero signed discrepancy. Zero observed counts
or equal rounded probability estimates do not establish that restriction.
The financial range is a continuation-value range, rather than merely the
current period's margin. An error that changes a default probability can
alter a terminal loss even when current margin is unchanged.

An exact three-period law illustrates the distinction. Begin in active state
$A$, with absorbing states $D$ and $P$ for default and prepayment. Both laws
earn USD 20 at each active opening and assign terminal rewards
$g(A)=0$, $g(D)=-1000$ and $g(P)=-50$. Their active transition rows are
\begin{equation}
 \begin{array}{c|ccc}
       & A & D & P\\\hline
 P_t(\cdot\mid A)&0.70&0.10&0.20\\
 Q_t(\cdot\mid A)&0.70&0.11&0.19
 \end{array}
 \qquad t=0,1,2.
 \label{exec17:s5-occupied-laws}
\end{equation}
The approximation moves probability $0.01$ from prepayment to default. It
preserves active survival exactly, so both laws have opening occupancies
$(1,0.7,0.49)$ and restricted active time $2.19$. Their period-three event
vectors $(F_D,F_P,S)$ are $(0.219,0.438,0.343)$ and
$(0.2409,0.4161,0.343)$. Their values are respectively $-197.1$ and
$-217.905$ USD. A correct survival curve alone therefore misses a model-value
error of $20.805$ USD per initial account.

The continuation values make the error traceable:
\begin{center}
\begin{tabular}{c r r r r}
\toprule
$t$ & $\mu^P_t(A)$ & $v^Q_{t+1}(A)$ & $b_t(A)$ & full next-state range\\
\midrule
0 & 1 & $-169.15$ & $9.5$ & $950$\\
1 & $0.7$ & $-99.5$ & $9.5$ & $950$\\
2 & $0.49$ & $0$ & $9.5$ & $1000$\\
\bottomrule
\end{tabular}
\end{center}
At every active row the residual is
$(-0.01)(-1000)+(0.01)(-50)=9.5$ USD. Equation
\eqref{exec17:s5-occupied-identity} gives
$9.5(1+0.7+0.49)=20.805$ USD. The full-range occupied envelope is
$0.01\{950+0.7(950)+0.49(1000)\}=21.05$ USD. If the equality of the
active probabilities is part of the admitted model class, the signed error
can involve only $D$ and $P$, whose continuation range is $950$ USD at every
step. The restricted-support envelope then equals the exact $20.805$ USD.

By comparison, the earlier uniform first-split envelope, using $\epsilon=0.01$,
a running range of USD 20 and terminal range of USD 1000, is
\begin{equation}
 20(0+0.01+0.0199)+1000(0.029701)=30.299\ \mathrm{USD}.
 \label{exec17:s5-occupied-comparison}
\end{equation}
Both bounds are valid in this construction. The tighter answer uses
information about occupancy and the destination of the error that the
uniform envelope did not require. It is not a universally better bound in
every process. Complete event-history enumeration, the backward recursion
and the separate geometric-duration calculation give the same reference
values here. These are exact authored laws, without a fitted external
estimator or observed bank transition audit.~\citep{exec17_s5_occupied}

\section{Which histories can validate the transition budget?}
\label{exec17:s5-transport}

A useful error budget must refer to the histories reached under the requested
policy. A held-out prediction sample often has a different history mix.
Even when its outcomes truthfully measure the desired transition, averaging
its errors under the wrong mix can hide financially important failures.
This distribution question remains after any causal identification argument
has supplied the relevant conditional transition law.

At a fixed time, let $\rho_t(x)$ be a validation input distribution and suppose
$\mu^P_t(x)>0$ implies $\rho_t(x)>0$. For the true signed continuation
residual $b_t$, write $w_t(x)=\mu^P_t(x)/\rho_t(x)$. Cauchy--Schwarz gives
\begin{equation}
 \begin{split}
 \left|\sum_x\mu^P_t(x)b_t(x)\right|
 &=\left|\sum_x\rho_t(x)w_t(x)b_t(x)\right|\\
 &\leq C_t\left\{\sum_x\rho_t(x)b_t(x)^2\right\}^{1/2},\qquad
 C_t^2=\sum_x\frac{\mu^P_t(x)^2}{\rho_t(x)}.
 \end{split}
 \label{exec17:s5-transport-bound}
\end{equation}
Thus a validated upper bound on the squared residual can be transported only
with an appropriate occupancy factor. The inequality concerns the true
conditional residual, not an arbitrary fitted prediction loss. Estimating
that residual or its squared norm with noisy transition observations requires
its own statistical control. A small in-sample squared loss cannot simply be
inserted as the right-hand side of a guaranteed interval.

Consider two history categories with validation probabilities $(0.99,0.01)$,
policy occupancies $(0.1,0.9)$ and residuals $(0,100)$ USD. Validation's mean
absolute residual is only USD 1, while the policy's signed value residual is
USD 90. The validation squared norm is $100$ USD$^2$, and
\begin{equation}
 C_t^2=\frac{0.1^2}{0.99}+\frac{0.9^2}{0.01}
      =\frac{8020}{99},\qquad
 C_t\|b_t\|_{2,\rho}=\sqrt{802000/99}\approx90.0056\ \mathrm{USD}.
 \label{exec17:s5-transport-example}
\end{equation}
The transported bound correctly reflects the nearly complete movement of
decision-relevant occupancy into the rare validation category. It does not
interpret the small unweighted USD 1 as a policy-value guarantee. Such a
residual can be generated by changing a transition to a state whose
continuation reward differs by USD 100; the numbers describe a possible
finite transition law, not merely an algebraic pair of incompatible scores.

If validation never visits the second category, the density ratio there is
undefined. No finite value of $C_t$ repairs that lack of support. With target
occupancy $0.9$ and an otherwise unrestricted continuation range of USD 1000,
the unseen category alone can contribute up to USD 900 in absolute error.
The missing information can be addressed by sampling that category, changing
the policy to avoid it, imposing and defending a transition restriction, or
retaining the resulting broad error set. The interface should expose this
choice. Increasing draws from the first category supplies none of the
missing second-category transition information.~\citep{exec17_s5_occupied}

This support failure should be distinguished from two other problems. First,
an observational transition conditional on a chosen action need not equal
the interventional transition needed by the policy if hidden confounding
remains. Reweighting the history mix cannot resolve that scientific failure.
Second, a measured account state can become available only after the
servicing action. Conditioning on its corrected warehouse value can define
the wrong information history even with complete statistical support. A
validation design must establish transition meaning and the availability
clock before it establishes precision for a residual.

These distinctions suggest concrete uses for public and generated data. A
generated finite process can reveal all transition rows and policy
occupancies, so it can test both the exact identity and deliberate support
failures. A public longitudinal cohort can exercise history construction,
censoring, event order and reweighting diagnostics. If its assignment or
measurement process does not identify the requested causal rows, it can still
test those pipeline operations while leaving causal validation open. Public
data do not become bank transition truth through a change of column names.
Conversely, missing bank truth need not delay a useful semantic prototype
with a fully specified generated process.

\section{A finite validation interval and its longitudinal allocation}
\label{exec17:s5-validation}

The exact identity can support a prospective validation interval under
stronger, explicit sampling conditions. Fix $Q$ using data separate from
the validation sample, so its continuation values are fixed when validation
begins. Index the relevant time/state rows by $j=1,\ldots,m$. Suppose the
reference occupancies $w_j=\mu^P_{t_j}(x_j)$ are known, and obtain $n_j$
independent next-state draws $Y_{jk}$ from the specified reference row $P_j$.
The sample sizes are fixed before observing these draws. Independence
between different rows is not needed for the union bound below, but the
within-row bounded independent-draw premise is needed. Define
\begin{equation}
 \widehat b_j=\frac1{n_j}\sum_{k=1}^{n_j}v^Q_{t_j+1}(Y_{jk})
              -\sum_y Q_j(y)v^Q_{t_j+1}(y),\qquad
 \widehat B=\sum_jw_j\widehat b_j.
 \label{exec17:s5-validation-residual}
\end{equation}
Let $R_j$ bound the range of the sampled continuation value under the
reference row. The bounded-score concentration argument used earlier gives
\begin{center}
$\Pr(|\widehat b_j-b_j|>R_j\sqrt{\log(2m/\alpha)/(2n_j)})\leq\alpha/m$.
\end{center}
Apply the union bound over the rows and then the triangle inequality with
nonnegative occupancies. With probability at least $1-\alpha$,
\begin{equation}
 \left|\{V(P)-V(Q)\}-\widehat B\right|\leq
 H(\boldsymbol n):=
 \sqrt{\frac{\log(2m/\alpha)}2}
 \sum_{j=1}^m\frac{w_jR_j}{\sqrt{n_j}}.
 \label{exec17:s5-validation-interval}
\end{equation}
Under these premises $V(Q)+\widehat B\pm H$ is a finite reference-value
interval. The result is conditional on the separately trained model and the
known occupancies. It validates the model's value error from genuine
reference transitions; generating extra draws from $Q$ itself estimates a
different expectation and cannot validate $P-Q$.

Known occupancy is a material premise, not a harmless convenience. In the
three-period example it follows from the scientifically imposed active
probability $0.7$ under both laws. The validation class may allow the
default/prepayment allocation to differ by time, while retaining that known
active probability. One can then validate each active row separately without
assuming that its default fraction is constant across times. When occupancy
is estimated instead, its error multiplies the row residual. A valid analysis
needs a joint confidence region or another justified bound for these weights;
plugging estimated weights into Equation \eqref{exec17:s5-validation-interval}
does not preserve its stated guarantee automatically.

The independent draws can be obtained from a generated reference process,
using distinct trajectories or separate conditional-transition draws.
Observational subject histories generally do not provide arbitrary numbers
of independent transitions at chosen histories. Repeated rows from one
account, adaptive stopping after a favorable residual, estimated causal
weights and random survivor counts require corresponding sampling arguments.
They are not covered by declaring $n_j$ to be the number of available rows.
The public prototype can test the finite result in its declared experiment
before proposing the additional assumptions and costs of authentic
validation.

The interval also shows why equal validation effort at every time can be
inefficient. Put $a_j=w_jR_j$. If one row observation costs $c_j>0$ and the
available variable budget is $C$, consider the continuous planning problem
\begin{equation}
 \min_{n_j>0}\ \sum_j a_jn_j^{-1/2}
 \quad\text{subject to}\quad \sum_jc_jn_j=C.
 \label{exec17:s5-validation-design}
\end{equation}
For positive $a_j$ the objective is strictly convex. Its first-order
conditions are $a_j/(2n_j^{3/2})=\lambda c_j$. Solving them and using the
budget equality gives
\begin{equation}
 S=\sum_j a_j^{2/3}c_j^{1/3},\qquad
 n_j^*=\frac{C(a_j/c_j)^{2/3}}S,\qquad
 H^*=\sqrt{\frac{\log(2m/\alpha)}{2C}}\,S^{3/2}.
 \label{exec17:s5-validation-allocation}
\end{equation}
Convexity establishes the global continuous minimum. The $2/3$ exponent
arises from adding row confidence radii, each proportional to $n_j^{-1/2}$.
It need not equal an allocation that minimizes the variance of a different
estimator. Zero-occupancy or zero-range rows require no observations for this
particular value bound, although other scientific and operating checks may
still need them. Minimum row counts, capacities and integer observations
add constraints; the continuous solution is then a planning reference.

In the example, use three active rows, $\alpha=0.05$,
$w=(1,0.7,0.49)$ and full continuation ranges $(950,950,1000)$ USD. With
equal row costs and budget $30000$, equal allocation gives $10000$ draws at
each time and a simultaneous halfwidth of approximately USD $32.5680$.
The continuous allocation is approximately $(12337.92,9726.89,7935.18)$.
The feasible integer choice $(12338,9727,7935)$ uses the same budget and
has halfwidth USD $32.1752$.

For costs $(1,2,4)$ and budget $60000$ in the stipulated cost units, equal
counts of $8571$ spend $59997$ and give halfwidth USD $35.1783$. The
continuous allocation is approximately $(19905.57,9885.99,5080.61)$.
A feasible rounding and local leftover-allocation rule yields
$(19908,9886,5080)$, spends $60000$, and gives halfwidth USD $31.4016$.
This integer rule is feasible; no claim of exact integer optimality follows
from the continuous proof. The cost units describe transition-validation
effort. They do not establish staff prices, implementation costs or bank
return on investment.~\citep{exec17_s5_occupied}

These are prospective radii from fixed counts, not intervals observed in a
new validation experiment. Even if the estimated center happened to equal
the exact USD $20.805$ discrepancy, a USD $31.4016$ halfwidth would not
certify its sign. A useful design must therefore relate the financial
precision requirement to the sampling cost before commissioning expensive
validation. Better use of event-support restrictions or a justified
homogeneous-row model may reduce that cost, while introducing additional
scientific premises. The calculation makes those tradeoffs visible without
pretending that a large nominal sample guarantees a useful decision.

Finally, the interval addresses one fixed policy, a fixed continuation model
and the specified reference transition experiment. Two policy values require
joint control over their rows and their common or separate reference data.
Choosing a policy or fitting $Q$ using the same validation observations
changes the experiment. Numerical approximation of $v^Q$, reward-model
error, unsupported tails, uncertain natural-action semantics and unmeasured
assignment truth also remain distinct. They can receive justified additional
envelopes or explicit unavailable components; none is already included in
the row-sampling radius. An integration that preserves these distinctions
can explain exactly what a public prototype established and which
information an eventual bank experiment would still need.


% SOURCE docs/release/2026-10-29/source/docs/writing-repair/2026-10-29/src/9ac4870b17074c5a.tex CHAR 55661


\section{From the policy law to the software comparison}
\label{exec17:s5-end}

Longitudinal causal software is comparable only after the decision clock,
event law, intervention meaning, support, uncertainty unit and monetary
objective are aligned. The exact constructions show where a plausible scalar
can reverse a decision: an unsupported tail, survivor composition, an ordered
competing-event conversion, a natural action that is not a redraw, a singular
cap, a cumulative-weight translation, or a simulation error presented as a
confidence interval. A product proposal can close these semantic and design
gaps with public/generated fixtures and bounded derivations now. It cannot
claim that those fixtures establish bank truth, general calibration, or final
adoption.

The source-specific studies that follow apply these distinctions to candidate
workers. The R and Python comparisons examine intervention simulation,
targeting, resampling and failure handling. The Java, C++ and other-language
comparisons then separate predictive survival components from a causal policy
estimator. Their different roles matter more than a shared output name.


% SOURCE docs/release/2026-10-29/source/docs/writing-repair/2026-10-29/src/c3ac7823f9d2fa96.tex CHAR 0
\section{Preserving longitudinal policies when reusing gfoRmula}
\label{exec20:gformula-scope}
A bank choosing repeated servicing actions must preserve the joint evolution of customer condition, action and exit. A simulator can return a smooth risk curve while answering a different intervention question from the one the bank intended. The gfoRmula source provides a useful concrete reference because the simulation, intervention and bootstrap steps are inspectable. The selected November 2025 revision declares version 1.1.2, while its changelog leaves that version's release date as \texttt{TBD}. The comparison therefore concerns an immutable source revision, not a claim about the latest released service. The inherited names gformula and gfoRmula refer here to the same project and do not count as two independent implementations.~\citep{exec20_gformula_reuse}

\subsection{The target precedes the simulation settings}
\label{exec20:gformula-target}
For a terminal outcome, let $L_0$ denote the information available before a first action, $A_0$ that action, $L_1$ the next permitted measurement and $A_1$ a second action. A deterministic policy $d=(d_0,d_1)$ chooses $A_0=d_0(L_0)$ and $A_1=d_1(L_0,A_0,L_1)$. Write $Y^d$ for the resulting terminal outcome and $\mu(d)=E(Y^d)$ for its population mean. With consistency, sequential exchangeability, positivity for the requested actions and an appropriate observation process, the earlier longitudinal argument gives
\begin{equation}
\begin{split}
 \mu(d)=\sum_{l_0,l_1}&m\{l_0,d_0(l_0),l_1,
                  d_1(l_0,d_0(l_0),l_1)\}\\
 &{}\times p\{l_1\mid l_0,d_0(l_0)\}p(l_0),
\end{split}
\label{exec20:gformula-law}
\end{equation}
where $m(l_0,a_0,l_1,a_1)=E(Y\mid l_0,a_0,l_1,a_1)$. For continuous variables the corresponding sums become integrals. This is an identification statement under assumptions about the data-generating experiment; fitting a conditional model does not supply those assumptions. In particular, a financial condition affected by an earlier action can be a confounder for a later action. Its distribution must evolve under the earlier policy rather than be held at its observed marginal distribution.~\citep{hernan2020}

The inspected entry point distinguishes survival risks from continuous and binary end-of-follow-up outcomes. These objects require different output types. The simulation takes baseline covariates from the observed baseline rows; at later times it generates covariates in the supplied \texttt{covnames} order, applies interventions, refreshes relevant histories and predicts outcomes. Thus a model specification includes the within-period ordering and history definitions, not merely a set of regression formulas. A treatment or balance field whose availability follows the outcome decision is not made admissible by sorting the rows. Custom outcome prediction must return the required probability or conditional mean; accepting a flexible prediction callback establishes neither causal identification nor interval calibration.~\citep{exec20_gformula_reuse}

Three sample sizes should remain explicit. The number $n$ of observed people governs the information available for fitting. The \texttt{nsimul} value $M$ governs the simulated population used in numerical integration. The \texttt{nsamples} value $B$ governs requested bootstrap replicates. Larger $M$ can reduce Monte Carlo error conditional on fitted models. Larger $B$ can approximate a chosen resampling distribution more closely. Neither change creates additional independent observed histories or remedies unsupported policy actions. The adapter should report all three quantities alongside the source population and time horizon.

\subsection{An exact two-period policy comparison}
\label{exec20:gformula-example}
Consider an authored binary law representing two opportunities for a helpful servicing action. Let $Y=1$ denote a beneficial terminal result. The variables and probabilities are stipulated, not estimated from a bank:
\begin{align}
 P(L_0=1)&=\tfrac12, \nonumber\\
 P(L_1=1\mid L_0,A_0)&=\tfrac15+\tfrac25L_0+\tfrac15A_0, \nonumber\\
 m(L_0,A_0,L_1,A_1)&=\tfrac1{20}+\tfrac15L_1
                   +\tfrac3{20}A_0+\tfrac1{10}A_1
                   +\tfrac3{10}L_1A_1.
 \label{exec20:gformula-example-law}
\end{align}
Every displayed probability lies in $[0,1]$ on the binary support. Define never acting by $A_0=A_1=0$, always acting by $A_0=A_1=1$, and a dynamic policy by $A_0=1-L_0$, $A_1=L_1$. In the never arm, $E(L_1)=2/5$, so $\mu_{\rm never}=1/20+(1/5)(2/5)=13/100$. In the always arm, $E(L_1)=3/5$, giving $\mu_{\rm always}=3/5$. For the dynamic arm, the two conditional probabilities of $L_1=1$ are $2/5$ and $3/5$; hence $E(L_1)=1/2$ and $E(A_0)=1/2$. Since $L_1A_1=L_1^2=L_1$ under this binary policy,
\begin{equation}
 \mu_{\rm dynamic}
 =\tfrac1{20}+\tfrac3{20}\tfrac12
    +\left(\tfrac15+\tfrac1{10}+\tfrac3{10}\right)\tfrac12
 =\tfrac{17}{40}.
 \label{exec20:gformula-values}
\end{equation}
These rational values also follow by enumerating every possible history. They are independent numerical references for a future simulator; they are not observed gfoRmula outputs.~\citep{exec20_gformula_reference}

A tempting shortcut replaces $L_1$ by its conditional mean $p$ and applies the equivalent binary-state rule $A_1=1\{L_1\geq1/2\}$ to that mean. At $L_0=0$, this gives $p=2/5$, $A_0=1$ and $A_1=0$, with predicted outcome probability $7/25$. At $L_0=1$, it gives $p=3/5$, $A_0=0$ and $A_1=1$, with probability $9/20$. Their average is $73/200$, rather than $17/40$. Integrating over individual histories and intervening on a mean history are different operations. The discrepancy follows from the thresholded action and the state--action interaction, without any simulation noise.

Suppose, solely for a decision illustration, that an additional beneficial result is worth USD 100 and assignment to this dynamic policy costs an additional fixed USD 20 per eligible account relative to never acting. The correct incremental margin is
\begin{equation}
 100\left(\tfrac{17}{40}-\tfrac{13}{100}\right)-20
   =\text{USD }9.50.
 \label{exec20:gformula-margin}
\end{equation}
The mean-history shortcut gives USD 3.50. Actual action-specific costs could require integrating a separate reward at each time; the fixed USD 20 assumption deliberately isolates the effect of preserving the outcome law. Neither figure is measured bank value.~\citep{exec20_gformula_reference}

\subsection{Competing events define a different decision}
\label{exec20:gformula-events}
Suppose default is the event of interest and an account's early repayment prevents a later default on that account. Early repayment also changes revenue and exposure, so treating it as ordinary loss of observation may change both the causal question and the economic calculation. In the inspected survival branch, the same-period main-event contribution is $p_{Y,t}(1-p_{D,t})$, where the competing event takes precedence. Previous main-event and competing-event survival multiply subsequent contributions. With constant probabilities $p_Y=1/5$ and $p_D=1/10$ for two periods, this convention gives
\begin{equation}
 R_{\rm competing}
 =\tfrac15\tfrac9{10}
  +\left(\tfrac45\tfrac9{10}\right)\tfrac15\tfrac9{10}
 =\tfrac{387}{1250}=0.3096.
 \label{exec20:gformula-competing}
\end{equation}
Suppressing the competing-event mechanism while keeping the stipulated main-event probabilities gives $R_{\rm eliminate}=1/5+(4/5)(1/5)=9/25=0.36$. The difference is $63/1250$, or USD 50.40 at an assumed loss of USD 1,000 per main event. This is a difference between declared event worlds, not a correction that should automatically be applied to either curve.~\citep{exec20_gformula_reuse,exec20_gformula_reference}

The inspected \texttt{compevent\_cens} branch disables the competing hazard in intervention simulation and retains separate handling for diagnostics. Its estimand therefore requires explicit interpretation. An intervention that eliminates early repayment may be scientifically or operationally implausible; identifying its consequences can require assumptions beyond those for observed cumulative incidence. Censoring due to a missing report is yet another mechanism. A library, CLI or MCP response must carry the target, event priority and observation convention with the curve. A field named risk cannot safely conceal those distinctions.

\subsection{Histories, thresholds and range clipping}
\label{exec20:gformula-history}
The helper checks expect consecutive numeric time rows within each person. A monthly history with an omitted month therefore needs an explicit measurement and missingness model; silently closing the gap changes the decision clock. Without pre-baseline rows, lag construction can use a baseline value or a zero/reference convention. A structural absence of earlier treatment is different from a missing earlier measurement. For an absorbing status such as permanent closure, neither choice should be inherited merely because it is a software default.~\citep{exec20_gformula_reuse}

The static intervention function indexes its supplied vector at absolute time $t+1$. Selected intervention times therefore do not imply that the vector is indexed by the number of interventions already performed. The newer intervention argument names identify both policy number and variable. A conversion should preserve that association and reject an ambiguous or too-short schedule before calling a worker. The threshold function directly clamps the current value to its bounds. For natural values $(-1,2,-1)$, a lower bound of zero gives $(0,2,0)$. A different rule that retains the preceding positive level would give $(0,2,2)$. The threshold routine supplies no such direct absorbing state; feedback through simulated histories may still affect later natural values. A bank asking for a persistent commitment needs to specify that commitment explicitly.~\citep{exec20_gformula_reuse,exec20_gformula_reference}

The default \texttt{sim\_trunc=TRUE} also matters. For selected continuous covariate types, the inspected code clamps simulated draws to empirical observed ranges. Clamping maps a draw outside a range to an endpoint; it is not rejection-resampling from a truncated distribution. An authored discrete witness makes the distinction transparent. Let $X$ take $-2,0,2$ with probabilities $1/4,1/2,1/4$. Then
\begin{equation}
 E(X^2)=2,\qquad
 E\{\min(1,\max(-1,X))^2\}=\tfrac12.
 \label{exec20:gformula-clipping}
\end{equation}
At an assumed quadratic loss coefficient of USD 100, the expected losses differ by USD 150. This witness is not a claim about the package's fitted normal distribution or a measured software defect. It demonstrates why a clipping option belongs in the scientific specification. A prototype should compare declared bounds, clipped mass and decision values under justified alternatives, and should not hide extrapolation problems by silently changing the simulated law.~\citep{exec20_gformula_reference}

\subsection{Bootstrap success is a scientific denominator}
\label{exec20:gformula-uncertainty}
The inspected bootstrap resamples whole person identifiers and their associated histories, refits models, and adjusts the simulated population. This preserves within-person records for each sampled individual. It does not by itself represent common bank shocks, shared calendar effects or errors in a common audit source. The appropriate independent unit depends on the observation experiment. A larger person-level bootstrap cannot manufacture independent calendars.~\citep{exec20_gformula_reuse}

The failure wrapper warns and returns missing results for failed replicates. The downstream standard-deviation and quantile calculations remove missing values. This is an observable source-level rule, not evidence that any particular native run has failed. Its implication for a proposed interface is nevertheless concrete: report requested, completed and failed replicate counts, failure reasons and the aggregation policy. An interval from the surviving replicates must not be described as though every requested fit contributed.

For a simple authored illustration, take four replicate estimates $(0,1,2,3)$. Their mean is $3/2$ and their sample variance is $5/3$. Retaining only $(0,1)$ changes the mean to $1/2$ and the sample variance to $1/2$. Selective computational survival can therefore change the resampling distribution. This example does not estimate actual failure probabilities or interval coverage for gfoRmula. It motivates a qualification experiment that retains every request and tests whether failure depends on difficult histories, sparse actions or unstable model fits.~\citep{exec20_gformula_reference}

Monte Carlo error within each fit, sampling uncertainty across fits, model misspecification and lack of causal identification remain distinct. A proposed study should vary $M$ while holding the observed data fixed to assess numerical integration, then repeat independent observed datasets under a known mechanism to assess the complete estimation and interval procedure. Natural-course agreement is a useful diagnostic for the fitted observed law, but it cannot verify unobserved counterfactuals or exchangeability. The selected upstream tests mainly check that example calls do not raise errors; they have not been run here and do not replace an independent causal or numerical oracle.~\citep{exec20_gformula_reuse}

\subsection{A reference worker and a public qualification route}
\label{exec20:gformula-reuse}
The proposal retains the pinned original R implementation as a reference candidate. A future process-separated worker could preserve the existing intervention, history, visit and bootstrap behavior more directly than a new port. Direct in-process adoption would add runtime and dependency commitments; a port would duplicate scientifically material semantics and need independent reconciliation. Neither has demonstrated an advantage here. Rejecting the entire package would also be unjustified: the inspected methods are relevant to bounded longitudinal questions even though unrestricted dependent-bank inference remains unsupported. The source declares GPL-3 and R dependencies, but a complete distribution, dependency and rights review is still required before packaging or porting it.~\citep{exec20_gformula_reuse}

The proposed request contract includes person identity, event and availability clocks, ordered covariate models, intervention schedules, history conventions, observation and competing-event targets, restrictions, $n$, $M$, $B$, seeds and range-clipping settings. Custom callbacks need immutable, reviewed identities. The output should identify a terminal mean, terminal probability or cumulative-risk curve; name its reference policy; preserve contrasts and numerical warnings; and separate sampling from simulation uncertainty. Invalid clocks, unsupported target coding, unqualified dependence and unacceptable replicate loss require explicit failure or unsupported status. The same typed object should survive library, CLI and MCP transport.

The first qualification experiment can use generated public mechanisms with the exact two-period values above, then extend to competing events, missing visits and alternative history conventions. A second experiment should isolate threshold and schedule semantics, clipping and person resampling, including intentionally failed fits. A third should test serialization, deterministic seed handling where promised and sequential versus parallel reproducibility at a frozen environment. Licensed public longitudinal data can then assess realistic support, censoring, fit diagnostics and computational burden, while retaining that real-data causal truth is usually unavailable. Simulation, refitting, conversion and analyst review costs should all be measured on matched accepted outputs. These experiments justify a prototype comparison; they do not confer native qualification, independent acceptance or bank adoption.


% SOURCE docs/release/2026-10-29/source/docs/writing-repair/2026-10-29/src/d2347a0ffab006bb.tex CHAR 0
\section{A longitudinal policy parameter must retain its probability law}
\label{exec23:lmtp-scope}

Sequential servicing software becomes useful only after the intervention is
specified as a rule on histories. The inspected \texttt{lmtp} revision
\texttt{f04a2b47f46debc515ce4ae778e05ebfde922c44}, dated 6 May 2026 and
declaring version 1.5.4, supplies ordered histories, shifted data, cross-fitting,
weights and longitudinal estimation interfaces. Its proposal disposition is a
reference choice with no permission to infer that an adapter is implemented.
The authenticated package declaration and license are AGPL-3; the README's
GPL badge differs. In particular, a separate process is not a sufficient
argument about distribution or network-source obligations. A concrete proposed
arrangement and its dependency rights still require review.
\citep{exec22_lmtp_reuse}

For a binary action with natural probability $e_t(h)$ at history $h$, the
inspected \texttt{ipsi} function implements a randomized modification whose
induced probabilities take the form
\begin{equation}
 q^{\rm shift}_t(h;\delta)=
 \begin{cases}
 \delta e_t(h),&0\leq\delta\leq1,\\
 1-\{1-e_t(h)\}/\delta,&\delta>1.
 \end{cases}
 \label{exec23:lmtp-shift}
\end{equation}
This differs from the odds-tilt intervention
\begin{equation}
 q^{\rm odds}_t(h;\delta)=
 \frac{\delta e_t(h)}{1-e_t(h)+\delta e_t(h)}.
 \label{exec23:lmtp-odds}
\end{equation}
Both maps agree with the natural law at $\delta=1$. A parameter value alone
does not identify which policy is meant. The distinction matters even with
perfect numerical computation and no unmeasured confounding.
\citep{exec22_lmtp_reuse,exec23_policy_reference}

Consider an exact authored example with two binary servicing contacts,
$A_1,A_2$. At each step the natural draw is conditional on the already assigned
past, with $e_1=1/5$, $e_2(0)=1/5$ and $e_2(1)=3/5$. Independently randomize
the declared modification of that draw. There is no censoring, competing event
or latent confounding in this stipulated law. Gross value is
$20A_1+30A_2-10A_1A_2$ dollars per account, and contact cost is
$8(A_1+A_2)$. For $q_1=P(A_1=1)$ and $q_{2a}=P(A_2=1\mid A_1=a)$,
the expected net value is
\begin{equation}
 V(q)=12q_1+22\{(1-q_1)q_{20}+q_1q_{21}\}-10q_1q_{21}.
 \label{exec23:lmtp-value}
\end{equation}
It follows by expanding net value as $12A_1+22A_2-10A_1A_2$ and taking
the three corresponding moments. Enumerating all four trajectories gives an
independent arithmetic check.~\citep{exec23_policy_reference}

\begin{center}
\begin{tabular}{lrrrr}
\hline
Policy at $\delta=2$ & $q_1$ & $q_{20}$ & $q_{21}$ & Net USD/account\\
\hline
Natural law & $1/5$ & $1/5$ & $3/5$ & $184/25$\\
Risk-style shift & $3/5$ & $3/5$ & $4/5$ & $456/25$\\
Odds tilt & $1/3$ & $1/3$ & $3/4$ & $107/9$\\
\hline
\end{tabular}
\end{center}

Thus the natural and shifted values are USD 7.36 and 18.24, whereas the
odds-tilt value is approximately USD 11.89. Relative to the natural law the
increments are $272/25$ and $1019/225$ dollars. These are exact values of an
authored law, not package estimates or bank benefits. They demonstrate a
conversion error that an interface can prevent: replacing one named
intervention with another changes the target before uncertainty is considered.

\subsection{Histories, folds, numerical safeguards and uncertainty are distinct contracts}
\label{exec23:lmtp-conversion}

In the inspected interface, supplied shifted data overrides a shift function;
shifted censoring becomes one, and each shift function reads the original data
argument. A bridge must preserve the declared sequence and the role of those
inputs rather than silently reapply earlier modifications. Restricting the
history length $k$ changes which information nuisance regressions receive.
That is a computational choice requiring a scientific sufficiency argument;
it does not establish sequential exchangeability.
\citep{exec22_lmtp_reuse}

Source inspection finds a further distinction between intended and realized
cluster handling. Data construction writes \texttt{..i..lmtp\_id}, while an
outer-fold call reads \texttt{.\_lmtp\_id}. Inner learner cross-validation
checks and forwards its own IDs, and the final \texttt{ife} call receives
\texttt{task\$id}. These are different paths. The naming mismatch needs a
native reproduction before any observed consequence is asserted. It cannot
be repaired by pointing to the existence of IDs elsewhere. The fold helper
may lower the requested number of folds; at $V=1$ it uses the same training
and validation set. The adapter should retain realized folds and clusters
and qualify the sampling law rather than report only a requested option.

The documented prediction bound is $10^{-5}$, but the control function's
default is $10^5$. Following the estimator call sites changes the conclusion:
the inspected TMLE predictions pass a hard-coded $10^{-5}$, while the SDR
predictions pass a null bound. There is no basis here for claiming an observed
estimator output of 100,000. Density-ratio calculations cap a denominator at
0.999 and subsequently trim using a quantile over the supplied object.
Warning suppression, global trimming and cumulative ratio-weighted residuals
can affect diagnostics. None of these numerical devices proves support,
nuisance convergence rates or calibrated uncertainty.
\citep{exec22_lmtp_reuse}

The returned estimate is a policy value; its singular \texttt{estimate} field
must not be mistaken for an automatically formed contrast. A contrast between
two estimated policies also needs their covariance. Variance and contrast
semantics delegated to \texttt{ife}, full survival handling, learner internals
and parallel random-number behavior remain unqualified. Public known-law
histories can test policy mapping, fold integrity and contrast preservation
before expensive bank experiments. This is a concrete route to evaluating a
longitudinal reference while keeping adoption and coverage as later claims.


% SOURCE docs/release/2026-10-29/source/docs/writing-repair/2026-10-29/src/598a0cc19531126f.tex CHAR 0
\section{Longitudinal regime contrasts: preserving the effective uncertainty method}
\label{exec23:ltmle-regimes}

A servicing proposal can compare a fixed contact schedule with a rule that
responds to hardship observed after an earlier contact. The distinction matters
because the first contact may change both hardship and the second action.
The selected ltmle source provides an R reference for binary treatment regimes,
longitudinal treatment and censoring histories, and treatment-specific means.
Its \texttt{ltmle} interface handles a single regime or a pair; its
\texttt{ltmleMSM} interface also supports a working marginal structural model.
A general working model is a projection under specified weights and covariates;
it need not reproduce each separate regime mean. A saturated two-regime
parameterization is a special case. The number of policies alone does not
determine whether a working model is saturated.~\citep{exec23_ltmle_source}

The inspected revision declares version 1.3-0 and GPL-2. The main source,
package declaration and supporting files match the complete tree of commit
\texttt{338c029d}\allowbreak\texttt{ae9692ef}\allowbreak\texttt{20714125}\allowbreak\texttt{773da703}\allowbreak\texttt{7688993b}, dated 17 April 2023.
A new default-branch query restricted to the client cutoff returns this same
commit. Repository metadata also records a June 2026 push; that event alone
does not establish a newer default-branch estimator or equality with a CRAN
artifact. The declared imports are \texttt{Matrix} and \texttt{matrixStats};
\texttt{SuperLearner} is suggested and required by the inspected non-GLM path.
The default \texttt{SL.library="glm"} uses generalized linear models. It does
not request an ensemble merely because an ensemble option exists.
\citep{exec23_ltmle_source}

\subsection{A dynamic intervention and its information clock}

Let $H$ indicate baseline borrower risk, with $P(H=1)=1/3$. The first contact
is $A_0$, hardship measured after it is $L_1$, and the second contact is $A_1$.
The outcome $Y=1$ denotes default by the fixed horizon. The information order
is $H,A_0,L_1,A_1,C,Y$, where $C=1$ means the final outcome is observed.
For an exact authored mechanism, stipulate
\begin{align}
 P(L_1=1\mid H,A_0)
 &=\begin{cases}
 1/4,&(H,A_0)=(0,0),\\
 1/8,&(H,A_0)=(0,1),\\
 3/4,&(H,A_0)=(1,0),\\
 1/2,&(H,A_0)=(1,1),
 \end{cases}\label{exec23:ltmle-hardship}\\
 P(Y=1\mid H,A_0,L_1,A_1)
 &=\frac1{10}+\frac H{10}+\frac{L_1}{5}
       -\frac{A_0}{25}-\frac{A_1}{20}.
 \label{exec23:ltmle-default}
\end{align}
Every binary history has a valid default probability: the minimum is $1/100$
and the maximum is $2/5$. Conditional on the listed history, let $C$ be
independent of $Y$, with probability $9/10$ of observation when $L_1=0$ and
$3/4$ when $L_1=1$. This is a stipulated observation mechanism; absence of
outcome data is not redefined as absence of default.

Compare status quo $d^0=(0,0)$, targeted servicing
$d^T=(H,L_1)$, and broad contact $d^B=(1,1)$. The second targeted action reads
the hardship observed at its decision time, including any effect of the first
action on hardship. To estimate these targets from logged data, one would need
consistency, sequential exchangeability and treatment support, together with
the censoring assumptions and relevant estimator regularity. These assumptions
are separate from the mathematical truth of the stipulated mechanism.
Its regime mean is
\begin{equation}
 \theta_d=\sum_{h,l}P(H=h)P(L_1=l\mid h,A_0=d_0(h))
 E[Y\mid h,A_0=d_0(h),l,A_1=d_1(l)].
 \label{exec23:ltmle-gformula}
\end{equation}
The hardship means under the three policies are $5/12,1/3,1/4$.
Thus both path enumeration and the separate expectation identity
$\theta_d=1/10+1/30+E_d(L_1)/5-E_d(A_0)/25-E_d(A_1)/20$
give the same risks. If a default costs USD 1,000 and contacts cost USD 6
and USD 4 at the two times, respectively, define net value relative to status quo by
\begin{equation}
 V_d=1000(\theta_{d^0}-\theta_d)-6E_d(A_0)-4E_d(A_1).
 \label{exec23:ltmle-value}
\end{equation}
The default cost is illustrative and fixed; it excludes any unmodeled changes
in recovery, customer behavior or subsequent service cost.

\begin{center}\small
\begin{tabular}{lccc}
\toprule
Policy & Default risk & Contact cost (USD) & Net value (USD)\\
\midrule
Status quo & $13/60$ & $0$ & $0$\\
Targeted dynamic & $17/100$ & $10/3$ & $130/3$\\
Broad contact & $7/75$ & $10$ & $340/3$\\
\bottomrule
\end{tabular}\end{center}
These are exact per-account expectations under the assumed law, not estimates
returned by ltmle or evidence of bank benefit. The positive broad-contact value
depends on the assumed effect and cost parameters. The example gives a target
for a future package comparison; it does not select a bank policy.
\citep{exec23_ltmle_exact}

Censoring alters even this small example if the analyst drops missing outcomes.
The complete-case risk under status quo is $69/335$, compared with the target
$13/60$; its bias is $-43/4020$. Under targeted servicing, the corresponding
risks are $139/850$ and $17/100$, giving bias $-11/1700$. Write
$\pi(L_1)=P(C=1\mid L_1)$. Under the stipulated conditional independence,
\begin{equation}
 E_d\!\left\{\frac{CY}{\pi(L_1)}\right\}
 =E_d\!\left[E\!\left\{\frac{CY}{\pi(L_1)}\,
          \middle|\,H,A_0,L_1,A_1\right\}\right]
 =E_d(Y)=\theta_d.
 \label{exec23:ltmle-censoring}
\end{equation}
Exhaustive enumeration verifies this population identity for all three
policies and checks that their probability masses sum to one. Known
observation probabilities make this an algebraic check. Estimating those
probabilities, multiplying treatment probabilities, and obtaining useful
uncertainty require additional work. The bound $\pi\geq3/4$ in this example
does not demonstrate overlap in a real servicing log.
\citep{exec23_ltmle_exact}

\subsection{Why a named variance option is insufficient}

The source defaults to \texttt{variance.method="tmle"}, but
\texttt{CreateInputs} switches to \texttt{ic} when the additional variance
method is unavailable. In the inspected revision, non-null \texttt{id},
nonbinary outcomes, \texttt{gcomp}, stratification and a deterministic-Q
function each make that additional method unavailable. The fallback itself
does not always warn: a separate warning path checks selected positivity and
rare-event indicators. The absence of that warning is therefore not evidence
that the caller's requested method ran.~\citep{exec23_ltmle_source}

For a single-regime TMLE summary, the code takes the larger of the available
influence-curve and additional variance estimates. For paired regimes, it
applies each contrast gradient to the covariance matrices and selects the
larger available standard deviation for that reported measure. It does not
justify a new covariance matrix formed by taking every element's maximum.
Risk and odds ratios are removed for nonbinary outcomes. Transformations back
to the outcome scale, including the distinction between mean and difference,
must be preserved. A larger variance estimate alone is not a proof of
coverage.~\citep{exec23_ltmle_source}

For clustered records, \texttt{HouseholdIC} sums record contributions within
each identifier and multiplies them by $m/n$, where $m$ is the number of
identifiers and $n$ is the number of records. Subsequent summary calculations
use the number of influence rows. This keeps the record-weighted target;
it does not automatically change that target to an equally weighted mean
of households of different sizes. The choice of weights and scientific
sampling unit must precede the conversion.

A separate exact covariance example uses four household influence pairs:
$(-3,-2),(-1,0),(1,0),(3,2)$. Each pair appears in two records. These are
hypothetical influence values, not fitted ltmle output. Aggregation and
rescaling recover the four pairs. Dividing their sample covariance by four
gives the covariance of the two means as
\begin{equation}
 \widehat\Sigma=\begin{pmatrix}5/3&1\\1&2/3\end{pmatrix},\qquad
 \widehat{\operatorname{Var}}(\widehat\theta_1-\widehat\theta_0)
 =5/3+2/3-2=1/3.
 \label{exec23:ltmle-cluster}
\end{equation}
Dropping the covariance instead gives $7/3$. Treating the eight duplicated
records as independent gives variance $5/7$ for the first mean, instead of
$5/3$. Their exact ratio is $7/3$ because sample variance uses different
degrees of freedom. The often-quoted factor of two is the leading
population-variance comparison for two identical records per independent
household. Neither the four-household calculation nor a small-sample
$t$ critical value proves calibration. ltmle uses a $t$ critical value below
100 influence rows and a normal critical value otherwise; the sampling law
and nuisance remainder still need justification.
\citep{exec23_ltmle_source,exec23_ltmle_exact}

\subsection{A proposed R comparator and the alternatives}

A candidate interface should declare the node order, policy information set,
binary action coding, event/censoring coding, outcome range, target weights
and cluster identity. Dynamic rules must not read future outcomes merely
because the R function receives a complete row. Numeric censoring inputs
use zero for censored and one for uncensored; the named conversion helper
distinguishes that meaning from a censoring indicator. For absorbing events,
the source checks that later outcome nodes remain one before cleaning
post-event records. The shorthand that later data are ``ignored'' must not
erase that validation requirement.~\citep{exec23_ltmle_source}

The cumulative-probability floor defaults to $0.01$ after multiplying the
incremental factors. It is not a floor on each factor before multiplication.
The source returns bounded and unbounded cumulative probabilities, which
support diagnostics; neither operation proves positivity. The outcome-range
code can truncate values before scaling and has a division by the range
width, motivating an explicit positive-width contract for a future adapter.
No warning should be interpreted as a complete fit audit: inspected helpers
suppress some rank-deficient prediction and extreme-probability warnings,
and the targeting helper can suppress a nonconvergence warning before its
score-repair path. A failed matrix solve can return missing standard errors.
Unavailable uncertainty must remain an explicit result state.
\citep{exec23_ltmle_source}

The proposal retains this version as a technical reference and a candidate
native R comparator. Direct integration requires resolved dependencies,
distribution rights, observed fit and warning behavior, and matched tests
against the declared policy and covariance targets. The GPL-2 declaration
and license distinguish running software from distributing or modifying it;
a separate process is not, by itself, a legal compatibility determination.
An integration design must review the actual distribution and dependency
chain. Porting to Python, Java or C++ additionally requires parity for
history/rule construction, censoring, outcome scaling, cumulative bounds,
learner behavior, clustering and covariance. No measured deployment or
performance need currently justifies that larger porting cost.
\citep{exec23_ltmle_source}

lmtp and gfoRmula remain substantive alternatives with their own modified
treatment, simulation and uncertainty contracts. Their names do not establish
equivalent interventions or interchangeable output. Selection should follow
the requested regime and observation law, then compare full fitting,
diagnostics, uncertainty, failure handling and maintenance costs. This
source-level reference decision grants no native integration capability.
The exact examples specify what a proposed trial should preserve; package
execution and empirical qualification remain distinct evidence.



% SOURCE docs/release/2026-10-29/source/docs/writing-repair/2026-10-29/src/59c808408a5dd148.tex CHAR 0

\section{Longitudinal R workers: preserve the regime, event and sampling unit}
\label{sec:exec23-longr-workers}
A repeated collections decision is not a baseline treatment comparison with
extra rows. At each visit a bank observes an eligible history, chooses an
action, and subsequently sees repayment, default, prepayment or loss of
observation. The order matters: an alert generated after the action cannot
become its baseline confounder. Nor may prepayment become censoring merely
because both terminate the collection record. The relevant software comparison
therefore begins with the regime and the event law, then follows each candidate
through its actual source interfaces. Six inherited lineages are reconciled
here: grf, lmtp, ltmle, gfoRmula, stremr and survtmle. This is a bounded comparison,
not a complete census or a finding of native numerical parity.

Let $H_t$ denote information available immediately before action $A_t$; let
$g_t(a\mid H_t)$ be the proposed intervention distribution. A cumulative
incidence target is $F_D^g(K)=\Pr(T^g\leq K,J^g=D)$, where $D$ denotes default,
$J$ the first terminal event, and $K$ a specified horizon. Identification from
observations additionally needs consistency, appropriate exchangeability,
support and a justified observation mechanism. A software argument for
$g$, a fitted survival curve or a finite influence vector does not establish
these premises. Dynamic treatment, a modified natural action, baseline
treatment and intervention on monitoring are different inputs, even if their
final output is a probability.

\subsection{Six lineages and their qualified roles}
The existing grf reference choice is deliberately narrow: a proposed
all-population AIPW contrast under a binary, independent-customer profile.
It excludes survival, repeated decisions and general calendar dependence.
Its inclusion in an inherited candidate register does not qualify that
initial profile or transfer it to this survival question. A separate
survival-source and qualification study remains necessary. \citep{exec21_grf_reuse}

The lmtp source revision declares 1.5.4. Its policy interface makes natural
and shifted treatment histories and the reference policy central to the
target. Its reviewed controls also separate outer folds, learner folds and
influence-function identifiers. The existing finding about defaults followed
actual TMLE and SDR forwarding; it was not an observed failure at a fictitious
parameter value. The remaining survival delegates, the variance dependency
and the effect of identifier conventions require their own inspection and
registered native cases. AGPL terms and direct source inspection do not
constitute complete distribution clearance. \citep{exec22_lmtp_reuse}

The pinned ltmle 1.3-0 route supplies a different longitudinal interface,
including binary treatment regimes. Requested variance, effective variance
and warnings must be preserved separately; choosing a cluster field is not
proof of an arbitrary dependence theorem. A future worker should return its
node order, intervention matrix, actual variance route and warning record
along with estimates. The previous source and exact-reference review remains
the basis for this scoped comparison. \citep{exec23_ltmle_source,exec23_ltmle_exact}

The gfoRmula source declares 1.1.2 while its NEWS date remains to be determined.
Its covariate simulation order, intervention functions, competing-event
convention and empirical-range truncation all affect the generated law.
Observed sample size, simulation size and bootstrap count are three different
precision controls. Resampling whole people preserves repeated measurements
inside a person; it does not reproduce arbitrary common branch or calendar
shocks. Failed bootstrap replicates and their omission from interval
aggregation must remain visible. These findings favor a narrowly specified
original R worker over a premature rewrite. \citep{exec20_gformula_reuse,exec20_gformula_reference}

The two additional source inspections sharpen the choice. stremr 0.4 has
separate treatment and monitoring intervention probabilities and sequential
regression machinery. Its TMLE and iterative-TMLE flags cannot both be set.
Internally generated inverse weights pass through stabilization and truncation
branches; user-supplied inverse weights take a different path after intervention
attribute checks. An adapter must record which path was used and the effective
weights, rather than treating a requested truncation value as evidence that
every supplied weight was truncated. Optional weights are rejected inside
the TMLE/iterative branch; this does not establish the same restriction for
every g-computation route. Documentation's warning about \texttt{Qstratify}
and an existing source branch are a qualification question, not proof of a
universal runtime failure. \citep{exec23_longr_stremr}

The survtmle 1.1.1 public wrapper accepts one treatment vector, failure times,
event labels and baseline adjustment variables. It dispatches to hazard or
iterated-mean estimators and forwards \texttt{maxIter} to the hazard route,
not to the mean call. That interface must not silently become a sequence of
repeated interventions. Missing inputs, nonpositive times and more than two
observed treatment values are rejected by the inspected checks. Requested
treatment/event combinations require observed endpoints. Requesting hazard
g-computation produces a warning and proceeds with TMLE, so requested and
effective methods can differ. The event-subset guard compares an expression
with itself; ordinary nonmissing values cannot activate that intended check.
This motivates a worker-side contract and native subset cases, not an
unobserved package-failure claim. \citep{exec23_longr_survtmle}

\subsection{Competing events change the monetary question}
\label{eq:exec23-longr-events}
Consider a supplied two-period law with mutually exclusive default and
prepayment hazards $d=1/10$ and $p=2/5$ at each period while still at risk.
The first-event default risk is
\[
 F_D(2)=d+(1-d-p)d=3/20.
\]
A different, stipulated law eliminates prepayment while holding the default
hazard fixed; its risk is
\[
 F_D^{\mathrm{elim}}(2)=d+(1-d)d=19/100.
\]
The second law does not follow from relabeling a prepayment field as censored.
It is an extra counterfactual model whose unchanged default mechanism needs
justification. Under a supplied USD 10,000 loss and a USD 1,700 contract
paying that loss on default, expected payout less price is minus USD 200
under the first law and plus USD 200 under the second. The sign reversal
comes from changing the target, not improving an estimator. These are exact
authored calculations, not observed banking economics or package output.
\citep{exec23_longr_reference}

For this reason, a prospective comparison must state whether it estimates
cumulative incidence with competing events, survival from all terminal
events, a cause-specific hazard, or a risk under an elimination intervention.
It must also specify within-period event ordering. The existing gfoRmula
inspection records its same-period competing-survival convention, whereas
the survtmle source builds cause-specific quantities from event labels.
A numerical match is meaningful only after these targets and clocks agree.
\citep{exec20_gformula_reuse,exec23_longr_survtmle}

\subsection{Stopping, warnings and residual custody}
A maximum-iteration option is not a convergence certificate. The inspected
stremr iterative path loops over $1,\ldots,m+1$. Its cap check is the
\texttt{else if} of the condition that the iteration exceeds one. For integer
$m\geq1$, persistent failure of its stopping test can therefore permit
$m+1$ updates. The influence contribution is evaluated before an update and
stored after the loop. At an exhausted cap, a worker must verify its relation
to the final predictions, rather than assume it was recomputed. Initial
success flags and failure counts are assembled before the optional iterative
fit; the later prediction update does not recompute those flags in that block.
\citep{exec23_longr_stremr}

The survtmle hazard loop checks $c\leq m$ and increments $c$ inside the body.
Persistent residuals can likewise permit $m+1$ updates, with an exhaustion
warning. Its fluctuation delegate has a separate optimizer cap. On nonzero
optimizer convergence it warns and substitutes zero coefficients before
continuing; the outer residual and the inner convergence status must both
remain visible. Other convergence checks in the hazard route can return
character sentinels, whereas the public wrapper subsequently expects list
fields. This source-level mismatch is a failure path to exercise; it is not
evidence that a specified native input has already reached it.
\citep{exec23_longr_survtmle}

An abstract loop witness with failed stopping and cap two produces three
updates for both control structures. A deliberately simple state marker
changes from two to three in the final update while a pre-update marker
remains two. This demonstrates the ordering issue only; it does not estimate
the magnitude of a TMLE error. A future worker should return requested cap,
actual updates, stop reason, final independently recomputed score and every
inner-optimizer warning. Changing upstream source would create a separately
pinned fork requiring a new audit. \citep{exec23_longr_reference}

The hazard influence code adds one column per requested treatment/event pair.
Its first extraction distinguishes numeric vectors from matrices, but its
post-update branch unconditionally calls \texttt{colMeans}. A single selected
column in an ordinary data frame is therefore a shape-sensitive path to
qualify. Initial identifiers are sequential; the long-data constructor expands
rows by subject, and influence contributions are summed by identifier. This
trace does not establish a row-permutation failure. Bounds joins, reused
models and later transformations still require explicit identity checks.
\citep{exec23_longr_survtmle}

\subsection{The unit of uncertainty survives reshaping}
\label{eq:exec23-longr-variance}
Let $D_{it}$ be a visit contribution and $D_i=\sum_t D_{it}$ a subject
contribution. A common influence variance expression is
\[
 \widehat V_{\mathrm{subj}}=n^{-2}\sum_i D_iD_i^\top.
\]
The inspected stremr and survtmle routes aggregate contributions within
subjects before outer products; their precise helpers and estimator branches
remain distinct. This supports preserving subject identity through the
interface. It does not justify deleting cross-subject covariance when subjects
share a branch shock. \citep{exec23_longr_stremr,exec23_longr_survtmle}

An exact law makes the distinction visible. Let independent signs $Z_1,Z_2$
take values $-1,1$ with equal probability. There are two subjects per branch
and two visits per subject, each visit contributing its branch sign. Dividing
the sum of the eight visit contributions by four subjects gives
$\widehat\theta=Z_1+Z_2$, with variance two. Treating eight visits as independent
while retaining that denominator gives $8/16=1/2$; aggregating by subject but
treating four subjects as independent gives $16/16=1$. Correct within-person
aggregation repairs one error while leaving common-branch covariance missing.
These values are exact under the supplied law, not empirical coverage results.
\citep{exec23_longr_reference}

Residual centering is another separate issue. With scalar contributions
$D_1=3,D_2=-1$, the uncentered sum divided by $n^2$ is $5/2$; the centered
version with the same denominator is two. Their difference is
$\bar D^2/n=1/2$. More generally the matrix difference is
$\bar D\bar D^\top/n$. This identity neither endorses an unbiased finite-sample
estimator nor repairs a failed estimating equation. It explains why returning
\texttt{meanIC}, final residuals and the exact covariance convention is more
informative than returning a standard error alone. \citep{exec23_longr_reference}

The inspected survtmle long-data constructor also deserves a clock contract.
Its observed expansion uses ranked unique failure times; its prediction
expansion uses the integer grid $1,\ldots,t_0$. With observed times two and
five, a late subject has two observed event-grid rows while the prediction
frame has five horizon rows. This is a difference in constructions, not a
demonstrated estimator error. Mapping irregular banking timestamps to time
indices, handling gaps and specifying the horizon must precede a parity trial.
The no-observed-censoring branch sets a censoring factor to one; absence of a
censoring label in an extract does not establish complete follow-up.
\citep{exec23_longr_survtmle,exec23_longr_reference}

\subsection{Integration choice and a bounded qualification program}
Retain the original R workers as candidate engines for their declared regimes.
A port to Python, Java or C++ would inherit the event ordering, nuisance fits,
weighting, failure behavior and uncertainty obligations; a shared API cannot
remove them. The initial preference is therefore isolated, pinned R workers
behind typed artifacts and the common CLI/MCP contract. A native port becomes
reasonable only when measured constraints justify it and an independent
reference covers the intended regime. The exact examples here serve as
semantic oracles, not as substitutes for upstream numerical conformance.

A request must bind the eligible cohort, time zero, subject and dependence
identifiers, available histories, intervention policy, competing events,
censoring, horizon and monetary units. It must additionally bind actual fold
membership, learner and dependency versions, weights, truncation, seeds,
variance method and resource limits. Responses preserve requested and effective
methods, row masks, effective weights, convergence status, failed bootstrap
counts and covariance. Unsupported regimes and unavailable engines must return
explicit refusals; an agent must not reinterpret them as lack of causal
identification or silently try a different estimand.

Qualification begins with a declared finite law containing treatment-confounder
feedback, monitoring, censoring and competing events, with target values
computed independently. Subsequent cases vary support, event subsets, single
output shapes, warnings, exhausted caps, duplicate identities, irregular
clocks and clustered shocks. A held-out simulation study should freeze
replications, seeds, acceptance thresholds and failure denominators before
running a locked environment. End-to-end costs include conversion, fitting,
bootstrap failures, validation, cancellation and restart. None has been
measured merely by inspecting source.

Rights and coverage are independent prerequisites. stremr declares GPL version
two or later and compilation through Rcpp; survtmle declares MIT plus a license
file containing year and copyright-holder fields. Those observations do not
clear the complete dependency distribution. Nor do the old selected CRAN
mirror dates prove abandonment. The saved discovery search counted 151 results
and returned ten; its unreturned tail and the named continuous-time lead
remain explicit survey follow-ups. Original scholarly texts must be acquired
and inspected before their technical results enter this comparison.
The bounded source comparison is ready for proposal use; exhaustive coverage,
native execution, calibrated uncertainty and bank adoption remain open.
\citep{exec23_longr_stremr,exec23_longr_survtmle}


% SOURCE docs/release/2026-10-29/source/docs/writing-repair/2026-10-29/src/0cdefa558bd74864.tex CHAR 0
\section{Python longitudinal inference: clocks, competing events and reproducible refusals}
\label{sec:exec24-longitudinal-python}
Repeated collections decisions require more than a survival curve. The bank
must specify when eligibility begins, which information is available before
each action, whether prepayment ends exposure to default, and which borrowers
remain in the population being averaged. The Python implementations examined
here offer different computational routes through that problem. zEpid 0.9.1
provides distinct baseline-survival, Monte Carlo longitudinal and iterative
conditional interfaces at its selected cutoff commit. pygformula 1.2.1
provides a parametric longitudinal simulator with intervention, competing-event
and bootstrap machinery. Their source contracts support a bounded comparison;
they do not establish complete Python coverage or native conformance.
\citep{exec24_longpy_zepid_source,exec24_longpy_pygformula_source}

\subsection{A causal clock and an event ordering}
Write $L_t$ for information observed just before decision $A_t$, and let
$H_t=(\bar L_t,\bar A_{t-1})$ denote the available history. A specified policy
sets $A_t=g_t(H_t)$. Sequential exchangeability, consistency, suitable support
and correctly specified transition models are substantive requirements for a
parametric g-formula interpretation; the simulator's ability to assign $A_t$
does not verify them. Censoring and competing events require their own
observation and intervention interpretations.

For a transparent authored law, suppose each period first permits prepayment
with conditional probability $q_t$, then default with probability $p_t$ among
those who did not prepay. With constant $q=1/5$ and $p=1/10$, the probability
of surviving both events for one period is $(1-q)(1-p)=18/25$. The two-period
default cumulative incidence is therefore
\[
 F_D(2)=(1-q)p+(1-q)(1-p)(1-q)p
       =\frac{86}{625}=0.1376.
\]
An intervention eliminating prepayment, while keeping the stipulated default
hazard unchanged, gives instead $1-(1-p)^2=19/100$. Under a stipulated loss
of USD 10,000 and protection premium of USD 1,600, the expected avoided-loss
minus premium is USD $-224$ for the first target and USD 300 for the second.
The latter describes a different hypothetical world; it is not the bank's
observed default risk. These exact invented examples establish neither real
insurance value nor identification from an observed bank panel.
\citep{exec24_longpy_reference}

The inspected pygformula simulator implements a conditional event ordering
that must remain visible. Its competing-event risk sums the default hazard
multiplied by survival of earlier default and competing events and by
survival of the current competing event. Without that branch it sums the
ordinary default-hazard increments. The constructor disables the explicit
competing-event branch when the option to treat competing events as censoring
is selected. That switch cannot be presented as an interchangeable numerical
option: an interpretation as elimination requires appropriate extra causal
and observation assumptions. Baseline time zero, pre-baseline histories and
lag initialization must also be preserved rather than inferred from calendar
strings.~\citep{exec24_longpy_pygformula_source}

\subsection{zEpid separates three estimation interfaces}
At commit \texttt{16a0f96f8b2c65df8715085801f21757d1478e1e}, the version file
identifies zEpid 0.9.1. The source date is a revision observation, not proof
of abandonment or a support guarantee. Its \texttt{SurvivalGFormula} describes
a baseline-fixed binary exposure and pooled logistic outcome model on long
person-period data. The fit overwrites exposure under the requested strategy,
predicts one-period survival, takes within-person cumulative products and
averages cumulative risks over the rows present at each time. Its documentation
requests an external bootstrap for confidence intervals. This interface is
not automatically a repeated-action policy estimator merely because its
input has repeated rows.~\citep{exec24_longpy_zepid_source}

The averaging population deserves a separate check. The inspected input
helper drops rows with missing values in any supplied column on this path,
including columns irrelevant to a chosen formula. A baseline-standardized
curve can consequently differ from a mean over the remaining records at a
later date. Consider equal baseline strata with supplied constant hazards
$1/10$ and $2/5$. Their two-period risks are $19/100$ and $16/25$, so the
baseline-standardized risk is $83/200=0.415$. If only the first stratum has
a retained row at the horizon, averaging the available risk returns $0.19$,
a change of $9/40=0.225$ due solely to the averaging population. This is an
exact composition fixture, not an observed output from a fitted zEpid model.
A proposed adapter must declare the standardization population, retain the
original identifiers and exclusions, and either build the justified prediction
grid or refuse a baseline-population label. Merely filling rows does not
resolve informative censoring.~\citep{exec24_longpy_zepid_source,exec24_longpy_reference}

The \texttt{MonteCarloGFormula} interface instead samples initial subject rows
and evolves ordered covariate models, exposure, outcome, censoring and lags.
Its source validates binary exposure and outcome; a competing-event branch
is a placeholder unreachable through the inspected constructor's nonbinary
outcome rejection. The initial sampled rows can be weighted, and the optional
low-memory output retains terminal rows rather than every simulated step.
Without a censoring model the simulation keeps records uncensored until its
administrative end. With a fitted censoring model, it draws the censoring
indicator after the outcome draw and sets that step's outcome to zero when
censoring occurs. These are precise ordering and target questions for a
matched trial, not evidence that every missing-outcome interpretation is
supported.~\citep{exec24_longpy_zepid_source}

Custom strategies and recodes in these interfaces use expressions evaluated
against the data frame, including explicit expression evaluation and code
execution hooks. In a proposed CLI or MCP workflow, a business-language agent
must not send arbitrary strings into those hooks. The proposal should accept
a typed policy with approved variables, operators, time availability and
registered transformations, then use a reviewed worker to construct the
permitted call. The contract must distinguish natural-course simulation from
retaining each observed action, bind random-number state and record all
intervention changes. Neither an executable expression nor an agent's
plausible explanation proves sequential exchangeability or positivity.
\citep{exec24_longpy_zepid_source}

The \texttt{IterativeCondGFormula} provides a further, wide-data route with
reverse sequential regressions. It rejects recurrent outcomes under its
binary-event check. The inspected two-dimensional treatment-plan validation
compares the plan's second dimension with the number of subjects, then also
requires that dimension to equal the number of exposure variables. For three
subjects and two decision times, an ordinary $3\times2$ per-subject plan
fails the first predicate. A one-dimensional plan is expanded separately.
The source also constructs a default-index data frame for the treatment plan
before assignment and iterates using a zip of exposures, outcomes, model
formulas and plan columns. Qualification should therefore check nondefault
indexes and unequal formula counts as well as shape. The arithmetic fixture
checks the stated predicate only; this review does not execute the package
or report a measured runtime failure.~\citep{exec24_longpy_zepid_source,exec24_longpy_reference}

\subsection{pygformula: simulation is only one uncertainty layer}
The selected pygformula commit is
\texttt{e66f787fcef2d6a9092b88135147c6b13785795c}.
Its constructor distinguishes survival, binary end-of-follow-up and continuous
end-of-follow-up outcomes, histories, observation visits and interventions.
The source defaults the horizon to the maximum time index plus one, inserts
a natural-course comparator, and constructs simulation and bootstrap seeds.
The simulator applies interventions at their specified dates and updates
histories and modeled covariates. Predefined static and threshold interventions
and custom callable paths require different conversion checks; the word
policy should not hide their information clocks or computational behavior.
\citep{exec24_longpy_pygformula_source}

Simulation precision is not fitted-model uncertainty. Under the elementary
invented Bernoulli law with risk $r=86/625$, an independent average of $M$
draws has conditional variance $r(1-r)/M$. Increasing $M$ addresses that
conditional simulation layer. It does not remove error in fitted transition
models, identify an unsupported counterfactual, or add independent branches
to the observed sample. The inspected bootstrap resamples subject identifiers
with replacement and assigns new identifiers to sampled copies before
refitting. That preserves within-subject histories; it does not by itself
qualify dependence across borrowers sharing a branch or calendar shock.
\citep{exec24_longpy_reference,exec24_longpy_pygformula_source}

Failure accounting is equally material. The bootstrap helper catches an
exception, warns with the replicate identity and returns missing result
fields. The caller first filters missing results into an accepted-result list,
but its subsequent estimate dictionary indexes every requested replicate
without the same guard. A failed replicate therefore needs an explicit
qualification case; it must not be silently turned into zero, counted as a
successful draw, or discarded without recording the denominator. Even a
simple designed pair of replicate values $(-1,1)$ has mean zero, whereas
retaining only $1$ has mean one. This example illustrates possible selection,
not the distribution of any package estimator. A production proposal needs
requested, attempted, accepted and failed counts, reasons and a predeclared
rule for withholding uncertainty.~\citep{exec24_longpy_pygformula_source,exec24_longpy_reference}

The source reuses a simulation seed across intervention comparisons, but
reuse of an initial seed alone does not establish a useful common-random-number
coupling when paths consume random draws differently. Native comparisons
must record the entire fitting, history, simulation and bootstrap environment;
include serial/parallel agreement and restart behavior; and distinguish
confidence intervals for a single target from simultaneous policy or time-path
claims. Natural-course agreement is a diagnostic for a fitted observational
process, not a proof of counterfactual validity.
\citep{exec24_longpy_pygformula_source}

\subsection{Integration choices and the remaining survey boundary}
The conditional route is a narrow original-Python worker for each declared
profile. A whole-library port is not justified by an unmeasured promise of
speed or simpler packaging. The baseline zEpid curve, longitudinal Monte Carlo,
iterative regression and pygformula competing-event simulator need distinct
method identifiers and refusal rules. R longitudinal alternatives discussed
earlier remain comparators only after matching target, time grid, observation
mechanism, sample weights, dependence and uncertainty; their existence does
not prove cross-language parity.~\citep{exec24_longpy_zepid_source,exec24_longpy_pygformula_source}

Both selected projects have inspected MIT notices, but that does not clear
all transitive software or model artifacts. zEpid's setup declares numerical,
data, regression and plotting dependencies with broad ranges. pygformula's
\texttt{pyproject.toml} declares version 1.2.1, Python at least 3.8 and a
setuptools build requirement, while its project dependencies include
\texttt{CompRiskReg}; the separate requirements file omits that dependency
and uses some more specific lower bounds. A lock cannot be synthesized by
assuming the two lists are identical. Native environment and distribution
review must resolve the actually selected installation route, all artifacts,
notices and supported numerical versions. Complete cost comparisons must
include conversion, fitting, failed fits, simulation, bootstrap, storage,
review and maintenance. No runtime advantage or bank ROI is established.
\citep{exec24_longpy_zepid_source,exec24_longpy_pygformula_source}

This bounded review leaves wider Python longitudinal and survival coverage
open. The saved discovery query returned only part of its counted records;
a refreshed cutoff-filtered query has a changed result set. Such metadata
provides leads, not technical evidence or a reconstruction of an earlier
index. Additional methods and identified originals require primary review,
and the selected implementations still need full helper/dependency inspection
and preregistered public/generated trials. Business-to-policy translation,
causal assumptions, numerical diagnostics and operating acceptance remain
separate responsibilities. The useful proposal result is a set of specific
integration and refusal contracts, with no claim that the whole area is closed.


% SOURCE docs/release/2026-10-29/source/docs/writing-repair/2026-10-29/src/06222592a872a00d.tex CHAR 0
\section{Java survival components: event clocks and defensible intervention services}
\label{sec:exec24-longitudinal-java}
Java survival components answer different parts of a banking problem. A
collections manager needs to compare action sequences for an eligible borrower
population; a capital analyst may first need event-time prediction. Combining
records held by several institutions introduces a confidentiality problem.
These requirements should not collapse into whether a package implements Cox
regression. Selected primary sources establish four candidates: H2O-3's
distributed fitting and portable scoring, BioJava's survival module,
FRESCO-stat's secure computation, and the distinct repository
\texttt{robbyjo/JLinAlg}. They do not establish exhaustive Java coverage,
qualified integration or a longitudinal causal estimator simply because a
survival routine exists.
\citep{exec24_longjava_h2o,exec24_longjava_biojava,exec24_longjava_fresco,exec24_longjava_jlinalg}

\subsection{Specify the event, clock and target first}
Let $L_i$ be entry time, $T_i$ the end of an observed interval, $D_i$ its
event indicator, and $X_i(t)$ information available before a decision at
$t$. In stratum $s$, stipulate the working model
$\lambda_i(t)=\lambda_{0s}(t)\exp\{X_i(t)^\top\beta+o_i(t)\}$, where
$o_i$ is a declared offset. Its risk set uses $L_i<t\leq T_i$. Neither
a fitted coefficient nor concordance identifies the consequence of changing
collections policy. A causal use still needs the treatment rule, population,
observation mechanism, sequential support and identification assumptions
developed earlier. A confounder affected by earlier treatment does not become
a valid ordinary adjustment variable merely because a start--stop row can
be supplied. The comparison separates the event-time computational component
from the intervention argument. The clock convention is explicit in the
selected H2O and JLinAlg source. \citep{exec24_longjava_h2o,exec24_longjava_jlinalg}

For a fixed baseline profile, write $r=\exp(X^\top\beta+o)$.
Integrating the stipulated hazard gives $S(t\mid X)=S_0(t)^r$.
A log score, relative score and probability have different units. With
$r=2$, baseline survivals $1/2$ and $3/4$ imply risks $3/4$ and $7/16$.
Under a stipulated USD~1,000 loss and USD~600 premium, margins are
$-150$ and $162.50$ dollars. The same relative score permits opposite
decisions. These are exact authored laws, not bank observations or fitted
package outputs. A CLI or MCP response must bind horizon, baseline, stratum
and centering before returning monetary consequences.
\citep{exec24_longjava_reference}

Event coding matters too. In a designed ten-borrower cohort, one default and
four prepayments give observed default fraction $1/10$; collapsing the events
gives $1/2$. Applying the same stipulated USD~1,000 default loss yields
USD~100 versus USD~500 per borrower. This coding witness is not a
counterfactual estimator. Prepayment as a competing event, censoring under an
observation model, and hypothetical elimination of prepayment remain distinct
targets. The earlier R and Python comparisons provide alternatives for
repeated-treatment requests. \citep{exec24_longjava_reference}

\subsection{H2O-3: separate fitting, baseline and scoring contracts}
The selected H2O commit is dated 20 September 2026; its build property says
3.47.0 and \texttt{bleeding\_edge}. This is a pinned development profile,
not a supported-release claim. Its Cox implementation accepts start and stop
columns, strata, offsets and observation weights, with Efron or Breslow ties.
The introductory document calls time columns integer-valued, but inspected
validation requires numeric columns and maps distinct stop times to indices.
The per-row discretizer rejects $L_i\geq T_i$; the preliminary global check
alone is not the full validation path. The common task skips bad or
zero-weight rows, while the processing routine rejects nonpositive weights
that reach it. A wrapper must report exclusions and resolved weights.
\citep{exec24_longjava_h2o}

The event routine subtracts the minimum observed event value and treats a
positive difference as an event. Multilevel codes are therefore not an
automatic competing-risk model. Require an explicit binary event mapping.
Coefficient covariance in the inspected path is the negative inverse Hessian;
the R summary bridge explicitly says robust variance was not used. Supplying
inverse-probability weights alone does not supply uncertainty for estimated
weights, repeated histories or branch shocks. \citep{exec24_longjava_h2o}

The documentation has a weighted-likelihood discrepancy. Its displayed
no-tie contribution is $\exp(w_i\eta_i)/\sum_{j\in R_i}w_j\exp(\eta_j)$,
whereas its Breslow expression and inspected Breslow code raise the denominator
to the event weight. With event weight two, another weight one, and relative
risks one and two, these expressions are $1/4$ and $1/16$. At covariates zero
and one with coefficient $\log 2$, the corresponding log-likelihood
derivatives are $-1/2$ and $-1$. This is not an ignorable constant.
Specify the implementation and tie/weight semantics, retain the discrepancy,
and require an independent fixture before native use.
\citep{exec24_longjava_h2o,exec24_longjava_reference}

Baseline outputs require another check. One cumulative-hazard array uses
tie-specific risk-set denominators. Named baseline frames use a separate loop
starting with total risk and subtracting rows at stop times; that loop does
not apply the Efron correction or add delayed entrants at entry times.
An abstract witness fixes unit risks: A fails at time one, B enters at one
and stops at two, and C is followed from zero to two. The declared risk set
at time one is A and C, giving increment $1/2$; an all-row initial total gives
$1/3$. Two tied events among three unit-risk rows separately give Breslow
$2/3$ and Efron $1/3+1/2=5/6$. These are source-path implications at
supplied risks, not observed H2O runtime failures. Test the exact exported
baseline rather than treating similarly named outputs as interchangeable.
\citep{exec24_longjava_h2o,exec24_longjava_reference}

The scoring frame is named \texttt{lp} and returns a centered linear
predictor. The MOJO scorer likewise returns a linear combination with
stratum centering and offset, not an absolute survival probability. Its
reader loads coefficients, means, strata and interactions; this score alone
does not provide a horizon-specific baseline. The source includes parity
tests for optional columns. One categorical/stratified test deliberately uses
Breslow because of a documented Java~11-plus Efron issue. That comment
identifies a qualification question, not failure in every current environment.
An original H2O fitting service and separately qualified portable scorer are
plausible choices; a whole-library port has no demonstrated benefit here.
\citep{exec24_longjava_h2o}

\subsection{BioJava: an existing port still needs semantic review}
The survival module belongs to a 7.3.1-SNAPSHOT parent. Source describes a
port of portions of R survival and explicitly says not all variants were
implemented. The selected helper accepts one event/censoring time, weights,
strata and clusters. Its row object stores one time; this path does not
establish a general counting-process interface. Controls expose ties,
scaling, tolerances and a nonconvergence flag. The helper skips nonpositive
weights before the downstream weight-use choice, catches exceptions, prints
a trace and returns null. A worker must preserve retained-row identities
and convert null into a typed failure, never a zero coefficient.
\citep{exec24_longjava_biojava}

The robust path saves naive variance, obtains score or dfbeta residuals,
groups by cluster when supplied and forms a cross-product. Strata conversion
still needs care: the fitting arrays describe terminal-row indicators,
whereas the helper reads integers and the comparator sorts by time and status.
Ordinary group labels require a verified mapping. Two branches also need
repair or refusal before reliance. A conjunction in the cluster guard can
dereference a null value before the intended error, while an empty string
passes. In robust unclustered cleanup, the loop saving predictors and scores
does not advance its index; restoration does. For abstract predictors $(2,3)$
and scores $(4,8)$, this structure restores $(3,0)$ and $(8,0)$. The cluster
branch does increment its index. This local trace does not claim all fitted
coefficients or robust variances are wrong: it identifies mutated row state
requiring a native fixture and reviewed repair.
\citep{exec24_longjava_biojava,exec24_longjava_reference}

Cluster identity is part of the sampling contract. For four scalar score
contributions $(1,1,-1,-1)$ and stipulated inverse-information factor $1/4$,
row independence gives $(1/4)^2(1+1+1+1)=1/4$. Grouping the first and last
pairs gives $(1/4)^2(2^2+(-2)^2)=1/2$. This is exact sandwich arithmetic,
not a coverage claim. Borrower clusters do not automatically cover branch
or calendar dependence. \citep{exec24_longjava_reference}

\subsection{FRESCO-stat: confidentiality is a separate contract}
At the pinned March~2023 commit, FRESCO-stat declares 0.6-SNAPSHOT,
FRESCO~1.3.3 and Java~9 source/target settings. Its survival row holds
secret-shared covariates, one time and a field named \texttt{censored}.
The gradient multiplies contributions by that field as an event-status bit;
the inspected example uses one for an observed event. Resolve the name
against behavior rather than guessing its encoding.
\citep{exec24_longjava_fresco}

The public class comment says gradient descent, but the optimizer computes
a Hessian pseudoinverse, multiplies it by the gradient, and subtracts a
scaled step for a fixed iteration count. The result has coefficients,
standard errors and a partial-likelihood-ratio statistic; the inspected loop
has no further convergence criterion. Fixed-point exponentials, inverses
and square roots add numerical obligations. The README warns that overflow
can produce very large outputs. A returned number does not certify
convergence or identified causal uncertainty. \citep{exec24_longjava_fresco}

The sorter identifies tied times using a helper that encrypts values under
a fresh shared key and opens ciphertexts. Equality groups then become public
values used in the risk-set calculation. Lists $(1,1,2)$ and $(10,10,20)$
have the same equality partition: their first two entries match. Revealing
that partition is different from reconstructing raw times, but the disclosure
contract must name it. This inspection proves no complete security theorem
or differential privacy for fitted outputs. Protocol suite, adversary,
modulus, precision, recipients and repeated-query policy remain required.
\citep{exec24_longjava_fresco,exec24_longjava_reference}

A separately qualified secure worker is a candidate when cross-custodian
analysis justifies its overhead. Centralized computation is a comparator
under a different access agreement. Secure computation does not create
treatment support, remove confounding or establish input truth. No upstream
protocol or test was executed in this comparison.
\citep{exec24_longjava_fresco}

\subsection{JLinAlg: explicit diagnostics and a narrower CLI}
The selected \texttt{robbyjo/JLinAlg} commit declares 0.3.6. Its vignette
distinguishes a general right-censored Cox/kinship-frailty CLI from Java APIs
for start--stop rows, strata and other models. The inspected fixed-effect API
validates finite $0\leq L_i<T_i$, nonnegative strata and at least one event.
It prepares risk-set structures and owns a backend in a closeable single-worker
fitter. The optimizer rescales predictors, uses step halving, tests coefficient
and score convergence, and returns convergence status and backend provenance
with covariance and baseline hazards. A result object need not have a true
convergence flag. \citep{exec24_longjava_jlinalg}

The counting-process baseline adds rows whose start is strictly before the
event and removes rows whose stop is before it; Efron and Breslow increments
are explicit. Separate diagnostics form cluster-summed score outer products
between covariance factors. They do not replace the ordinary errors already
in the fit result. The wrapper must name its selected covariance and cluster
population. The inspected fit signature has no case-weight argument; these
sources do not establish a general weighted marginal structural model.
Intervals and robust errors are components, not a sequential intervention
design. \citep{exec24_longjava_jlinalg}

The vignette reports development-host timings and broader frailty tests.
They are not independently reproduced here, and selected fixed-effect review
does not qualify mixed, sparse or GPU paths. The build pins a JDistlib~0.10.2
jar with an expected SHA-256, a useful artifact declaration that still does
not establish a complete locked environment or transitive rights. Backend
selection, helper numerics, runtime behavior and full costs remain open.
\citep{exec24_longjava_jlinalg}

\subsection{Adoption choices and executable next steps}
Direct texts declare Apache~2.0 for H2O, LGPL~2.1 in BioJava's root,
MIT for FRESCO-stat, and GPL~2 in JLinAlg's root with GPL-2.0-or-later
headers in the selected Java files. BioJava's POM says LGPL v2 while linking
the 2.1 text. Preserve these declarations and resolve distribution and
transitive artifacts with a rights owner; process separation is not legal
clearance. \citep{exec24_longjava_h2o,exec24_longjava_biojava,exec24_longjava_fresco,exec24_longjava_jlinalg}

The proposal choice is reference use of all four bounded profiles, followed
conditionally by original-language workers. H2O is a distributed-fitting and
portable-scoring candidate; BioJava offers inspectable robust machinery with
the stated branch and strata obligations; FRESCO-stat addresses a specialized
custodian problem; JLinAlg offers explicit interval and diagnostic APIs.
R/Python workers remain alternatives for repeated-treatment targets.
Reject automatic causal routing from a survival label, automatic log-score
conversion into money, and a whole-library port justified only by language
preference. \citep{exec24_longjava_h2o,exec24_longjava_biojava,exec24_longjava_fresco,exec24_longjava_jlinalg}

Execution starts with the method owner specifying target, clock, event
dictionary, observation model, treatment rule and dependence unit. The
adapter owner freezes a worker artifact and typed input/output schema,
including retained/rejected IDs, resolved options, convergence, covariance,
baseline, warnings and artifact identity. Agents may propose declarations;
unresolved semantics require a structured refusal before fitting. CLI jobs
and MCP calls should preserve the workflow chapter's request, cancellation
and artifact rules rather than silently converting failure into success.

The evaluation owner can preregister generated cases for delayed entry,
ties, nonunit weights, event coding, missing clusters, nonconvergence,
centering and covariance selection. Source-specific cases cover H2O's
baseline branches, BioJava's restoration, FRESCO's tie disclosure and numeric
range, and JLinAlg's separate diagnostics and backend closure. Compare with
an independently specified finite law or accepted reference implementation,
retaining failed requests in the denominator. Public data then support
conversion and realistic-size trials without pretending to supply known
causal truth. These exact witnesses motivate a prototype and its experiments;
they are not completed prototype evidence. \citep{exec24_longjava_reference}

Full costs include preparation, rights review, environment maintenance,
failed fits, diagnostics, protocol communication and operator time. Wider
search leads, uninspected delegates and native qualification remain explicit.
The comparison replaces an uninformative zero-result search with concrete
evidence and actions, without declaring the area closed.


% SOURCE docs/release/2026-10-29/source/docs/writing-repair/2026-10-29/src/2f545464f42acb7d.tex CHAR 0
\section{C++ survival components: preserving the event, clock and output}
\label{sec:exec24-longitudinal-cxx}

C++ survival components answer several different questions inside a banking
workflow. A collections team may want the observed probability of default by
month twelve, the default probability in a hypothetical world without prepayment,
or the effect of a sequence of contact decisions. These are different targets.
A fast predictor of recorded survival time does not choose among them. The
proposed product must preserve the event dictionary, observation and treatment
clocks, competing transitions, selection law and uncertainty target before it
chooses a computational engine. This comparison examines four pinned lineages:
XGBoost, CatBoost, ranger and aorsf. The evidence consists of authenticated primary
source and documentation, not execution of these packages or their tests.
Discovery remains bounded; neither four candidates nor a repository search
establishes exhaustive C++ coverage.~\citep{exec24_longcxx_xgboost,exec24_longcxx_catboost,exec24_longcxx_ranger,exec24_longcxx_aorsf}

\subsection{The estimand precedes the software label}

Let $T$ be the first transition time and $J$ its type, with $J=D$ for default
and $J=P$ for prepayment. The default cumulative incidence is
$F_D(t)=\Pr(T\leq t,J=D)$. It is not generally the complement of survival
obtained by treating prepayments as ordinary censoring. In an authored finite
cohort of ten accounts, four prepay at time one, two default at time two and
four remain open. With these complete observations, the default incidence is
$2/10$. The product of surviving all earlier transitions and the conditional
default fraction is $(6/10)(2/6)=1/5$. Treating the four prepayments as censored
instead gives the default-only complement $2/6=1/3$. This is a different
functional, not an arithmetic refinement of the same answer. At a stipulated
USD~1,000 default loss, the two expected-loss calculations are USD~200 and
USD~$1{,}000/3$. Neither a different language nor more trees resolves the
target substitution. A hypothetical intervention eliminating prepayment also
needs its own scientific assumptions; censoring a record is not such an
intervention.~\citep{exec24_longcxx_reference}

Observation intervals need equally careful interpretation. If account $i$
enters observation at $L_i$ and exits at $U_i$, the risk indicator for a
counting-process analysis is $Y_i(t)=1\{L_i<t\leq U_i\}$. For the three
authored intervals $(0,1]$, $(1,2]$ and $(0,2]$, with a default on the first
account at time one, the correct risk set contains two accounts. Its unit-risk
hazard increment is $1/2$, compared with $1/3$ if every terminal-time row is
treated as present from zero. Bounds on an unknown failure time are different
from entry and exit times: an interval-censored observation says that an event
occurred within a range, whereas delayed entry says when an account becomes
observable. An adapter must refuse their silent interchange. Repeated rows also
do not by themselves provide sequential exchangeability, positivity or
cluster-valid uncertainty for a treatment policy.~\citep{exec24_longcxx_reference}

\subsection{Boosted Cox and accelerated failure time objectives}

The selected XGBoost commit is
\texttt{c914bd00a946081c96e2f7aaa0e298f38b7b17e1}, dated 7 October 2026.
Its root build declares version 3.5.0; this identifies inspected development
bytes, not a supported binary release. The Cox objective sorts absolute labels,
uses positive labels for observed events and negative labels for censoring,
and exponentiates the margin for prediction. That exponentiated output is a
relative risk score. It does not contain a baseline survival curve or a monetary
default probability. The inspected routine multiplies each final gradient and
Hessian by its row weight, while its risk-set exponential sum is unweighted.
This source contract must not be relabelled as an inverse-probability weighted
longitudinal estimator without a separate derivation.~\citep{exec24_longcxx_xgboost}

XGBoost's AFT documentation instead specifies
$\log T=m(x)+\sigma Z$, with normal, logistic or extreme-value choices for
$Z$. Its ranged labels encode exact, right-, left- or interval-censored failure
times. The examples use zero as a left bound and positive infinity as the right
bound for right censoring. The inspected CPU and CUDA kernels call the same
distribution-specific derivative helpers and apply row weights. The prediction
transform returns $\exp(m(x))$, while the evaluation transform retains the raw
margin for the AFT loss. Declared CUDA kernels are evidence of implementation
text, not observed GPU conformance, speed or numerical parity. The pinned
tutorial routes ranged labels through the native training and data-matrix
interfaces and states a limitation for its scikit-learn interface. These
statements are tied to the acquired version.~\citep{exec24_longcxx_xgboost}

The AFT output requires its distribution and scale to be useful. Consider an
authored logistic AFT law with $\exp(m)=2$ and horizon eight. For scale one,
the failure probability is $F(\log 4)=4/5$; for scale two it is
$F(\log 2)=2/3$, where $F(z)=e^z/(1+e^z)$. Both models return the same
exponentiated location. A consumer that interprets that location as a probability
cannot distinguish them. The proposed result therefore carries the raw margin,
transformed location, distribution, scale, time unit and requested horizon as
separate fields.~\citep{exec24_longcxx_reference}

The CatBoost commit
\texttt{1ac99c589b54e91a349dd04a049f8ec1e1572326}, also dated 7 October
2026, declares Python package version 1.2.10. Its documentation presents Cox
and SurvivalAft objectives and marks their GPU support unavailable in the
inspected table. The displayed Cox expression uses $a_i$ inside an exponential,
and the AFT helper uses $(\log t-a)/\sigma$. These mathematical roles require
more precise output typing than the documentation's common phrase \emph{hazard
rates}. The raw objective margin is not an absolute hazard or probability.
CatBoost's inspected AFT routine uses $-1$ as an unbounded endpoint, with a
separate lower-end branch for left censoring. XGBoost's documented zero and
infinity conventions therefore cannot simply be copied into a CatBoost target
array. End-to-end target validation and the public prediction transformation
still need a frozen wrapper review.~\citep{exec24_longcxx_catboost}

Both projects clip or regularize AFT derivatives and provide limiting branches
for numerically extreme inputs. These choices are part of the computation,
not proofs of statistical calibration. CatBoost's selected Cox and AFT derivative
signatures leave their weight argument unused, but the inspected Cox caller
also passes weights into subsequent leaf accumulation. The appropriate finding
is that weighted-objective equivalence is unqualified: a local unused argument
does not establish that the entire training system ignores weights. A prototype
must compare the complete selected fitting path with an independently specified
weighted target.~\citep{exec24_longcxx_xgboost,exec24_longcxx_catboost}

\subsection{A tie fixture that can falsify silent equivalence}

Source review supplies a concrete acceptance question before a costly benchmark.
Take three unit risk scores, with two observed events tied at time one and a
third account censored at time two. For margins $f_i$, the Breslow negative
log-likelihood contribution is
\[
 \ell(f)=-f_1-f_2+2\log\{e^{f_1}+e^{f_2}+e^{f_3}\}.
\]
At zero margins its gradient is $(-1/3,-1/3,2/3)$. Its components sum to zero,
as required by invariance to a common additive margin. An independent rational
trace of the inspected XGBoost loop, in event/event/censor order, gives
$(-2/3,-1/3,2/3)$: the first tied event is processed before the second event's
reciprocal-risk contribution enters the running accumulator. The sum is
$-1/3$. The analogous inspected CatBoost loop removes the previous exponential
at each positive event and gives negative-log-likelihood-sign scores
$(-2/3,-1/6,5/6)$. Their sum is zero, but the two tied events are not symmetric
and the vector differs from the stated Breslow derivative. These are abstract
source-loop traces, not native benchmark results or a claim that every release,
wrapper or training regime behaves identically.~\citep{exec24_longcxx_xgboost,exec24_longcxx_catboost,exec24_longcxx_reference}

The next justified experiment is narrow: freeze the complete dispatch path,
evaluate equal-time permutations and event/censor ties, compare finite-difference
derivatives with an independent likelihood, and retain both the returned
gradients and row identities. Non-unit weights and extreme margins form separate
fixtures. A failing fixture would justify a restricted regime, another engine,
or a reviewed repair; it would not justify silently changing the scientific
likelihood. Until that experiment is admitted and performed, these findings
remain source-level concerns and acceptance criteria.

\subsection{Survival forests and transformations that change the decision}

ranger is pinned to
\texttt{dafa5db57972dd49c808db881e6e6cd5c4ebdc4a}, dated 16 January
2026; its R description declares 0.18.0. The inspected C++ leaf routine sums
$d_j/n_j$ to construct a cumulative hazard. The forest averages those hazards,
and the inspected R prediction wrapper returns survival as the exponential of
their negative. Its OOB error uses one minus concordance on summed cumulative
hazards for accounts with at least one OOB tree. That is a ranking diagnostic,
not an interval for a causal effect or a certification of horizon calibration.
The wrapper rejects nonbinary event coding in the selected survival path, and
its two standard-error choices explicitly reject survival forests. Case weights
are documented as sampling weights for bootstrap or subsampling selection.
These contracts are more specific than the labels \emph{weighted forest} or
\emph{uncertainty available}.~\citep{exec24_longcxx_ranger}

aorsf is pinned to
\texttt{085db7062d995eeafa4e210481d66af2eb02b8dc}, dated 5 February
2026; its description declares 0.1.6. Its documented accelerated survival
control uses a one-iteration Cox fit to construct split directions. The C++
delegate distinguishes that path from further iterations. This is a tree-building
choice, not a claim that the final forest is a converged Cox regression.
The selected leaf code maintains a product-limit survival estimate and an
additive cumulative-hazard estimate separately, with weighted event and risk
counts. The wrapper documents normalization of supplied weights to the number
of observations. Risk, survival, cumulative hazard, mortality and time are
distinct prediction types. Predictions beyond maximum follow-up are documented
as constant if boundary checks are disabled; a flat continuation is not evidence
about an unobserved future.~\citep{exec24_longcxx_aorsf}

For an authored leaf with one event among two at risk, the product-limit
survival is $1/2$, whereas the cumulative hazard is $1/2$ and its exponential
survival is $e^{-1/2}\simeq0.606531$. The corresponding risks are $0.5$ and
approximately $0.393469$. With a stipulated USD~450 premium and USD~1,000 loss,
the margins are USD~$-50$ and approximately USD~56.53. These are deliberately
simple constructed outputs; they are not observed errors or bank estimates.
They show why a consumer must retain the estimator and transformation, even
when both fields are called survival. Aggregation creates another distinction:
for two authored hazards $(0,\log4)$, exponentiating the negative average gives
$1/2$, while averaging their exponentiated negatives gives $5/8$. An adapter
must reproduce the declared aggregation order rather than substitute a familiar
formula.~\citep{exec24_longcxx_reference}

\subsection{Integration, adoption and justified rejection}

The proposed first route retains original engines behind typed workers. XGBoost
and CatBoost are conditional boosted-prediction alternatives once their objective,
endpoint, margin, weight, tie and environment contracts pass selected fixtures.
ranger's standalone C++ library or CLI and its R wrapper are separate packaging
choices; the inspected build declares C++14 and library/CLI targets. Its primary
headers explicitly distinguish the MIT C++ core from the GPL-3 R package.
XGBoost and CatBoost have Apache-2.0 root declarations. aorsf declares MIT and
R/Rcpp/RcppArmadillo dependencies. These direct declarations do not clear every
dependency, model artifact or redistribution. Rights review must bind the actual
chosen distribution.~\citep{exec24_longcxx_xgboost,exec24_longcxx_catboost,exec24_longcxx_ranger,exec24_longcxx_aorsf}

For aorsf, an original R worker preserves the formula, validation and prediction
semantics while using its C++ core. Extracting that core into a new standalone
port would require recreating the R-side contract and numeric dependencies;
the proposal has not demonstrated a benefit sufficient to justify that work.
For ranger, direct C++ embedding remains an alternative to an R worker when
deployment requirements and rights justify the separate interface. Whole-library
rewrites are rejected for the present proposal because they add parity and
maintenance obligations without demonstrated need. Native workers, a dedicated
longitudinal R/Python estimator and an independent finite-law reference remain
substantive alternatives, selected according to the target rather than language
preference.~\citep{exec24_longcxx_aorsf,exec24_longcxx_ranger}

The CLI and MCP proposal should expose a versioned request containing population,
time zero, availability clock, event dictionary, competing-event interpretation,
entry/censoring representation, treatment history, weights and their meaning,
sampling units, horizons, options and artifact hashes. The result must preserve
accepted and rejected row identities, objective and output type, baseline or
distribution parameters, convergence/numerical warnings, uncertainty method,
timing and complete request costs. Refusal is a typed result. A finite output,
successful process exit or replayed artifact must not be promoted to causal
identification or qualified uncertainty. These are proposed interfaces; no new
native adapter or sibling integration is implemented by this review.

The method owner must first decide whether a public or generated study can
identify the proposed target. The adapter owner then freezes version, wrapper
and conversion fixtures; the rights owner resolves the actual distribution;
the evaluation owner preregisters held-out targets, ties, censoring and competing
events, row permutations, dependent units, failure denominators and full costs.
No purchase, vendor trial or expensive bank experiment is necessary to write
those contracts. Candidate aliases, unreturned discovery tails, uninspected
delegates, cluster-calibrated uncertainty and native performance remain recorded
follow-ups. This bounded comparison improves the proposal's executable evidence
and actions without declaring C++ longitudinal coverage complete.


% SOURCE docs/release/2026-10-29/source/docs/writing-repair/2026-10-29/src/4576dc3f2d02826a.tex CHAR 0
\section{Julia and Rust survival components: preserving the clock and the estimand}
\label{sec:exec24-longitudinal-other}

A collections decision is defined by its clock before it is defined by its software language.
Consider accounts eligible for a payment holiday at the start of a month. The bank may
compare initiation within a grace period against continued ordinary servicing, follow
default and prepayment, and value interest receipts during a fixed horizon. A program
that returns a survival curve has solved only one part of this problem. The eligibility
date, treatment strategy, event definition, observation mechanism and sampling unit
must survive the translation from the business request to the estimator. This section
examines six concrete Julia and Rust alternatives at commits preceding the client
cutoff of 7 October 2026. Their role is conditional reuse or reference, not evidence
that a native causal product has already been qualified.

\subsection{A target and an observation contract}
Let $T^a$ denote event time under strategy $a$, measured from eligibility, and let
$S_a(t)=\Pr(T^a>t)$. Two possible targets are the event-risk contrast
\begin{equation}
 \Delta_F(\tau)=\Pr(T^1\leq\tau)-\Pr(T^0\leq\tau),
 \qquad
 \Delta_M(\tau)=\int_0^\tau\{S_1(t)-S_0(t)\}\,dt.
\end{equation}
The first is a probability difference; the second is a restricted-mean survival-time
difference and has time units. If the event is default, prepayment needs an explicit
competing-event or hypothetical-intervention interpretation. Treating prepayment as
ordinary independent censoring changes the question unless the necessary assumptions
are supplied. These definitions do not by themselves identify either counterfactual
curve. Consistency, treatment and censoring exchangeability under the specified
history, positivity and an appropriate observation law remain substantive conditions.

Write $L_i$ for entry, $C_i$ for right censoring and $\widetilde T_i$ for observed exit.
For a continuous clock with a declared boundary convention, a risk indicator might be
$Y_i(t)=\mathbf1\{L_i<t\leq\widetilde T_i\}$. By contrast, an interval-censored
record $(u_i,v_i]$ says that the event occurred somewhere in that interval; it does not
say that observation began at $u_i$. Three otherwise eligible accounts, one entering
after the first observed event, give a first-event increment $1/2$, while ignoring
entry gives $1/3$. This exact example concerns data meaning, not numerical speed.
The independent reference also constructs two deterministic event times, four and
two months, with horizon five. Their restricted means differ by two months while
their event risks at the horizon are both one, so the risk difference is zero.
Multiplying the time scale by thirty changes the mean contrast to sixty days and
leaves the probability contrast unchanged.~\citep{exec24_survivalother_reference}

The proposed command and MCP request must therefore bind population, eligibility
clock, intervention strategy, history availability, entry, censoring, competing events,
horizon, output units and sampling unit before selecting an engine. A convenient
two-column schema is insufficient evidence that two engines implement the same target.

\subsection{Julia alternatives and their distinct contracts}
The JuliaStats Survival.jl snapshot declares version 0.3.3 and Julia 1.10 compatibility,
with a direct MIT licence. Its exported event type contains a time and a Boolean
event/right-censoring indicator. The event-table constructors, Kaplan--Meier component
and Cox fitting component provide useful building blocks for a narrow right-censored
worker. They do not, through that response type alone, encode entry or a treatment
history. The inspected Cox objective uses fractional removal of tied event risk
scores, an Efron-style convention.~\citep{exec24_survivalother_juliastats}

Even this narrow candidate needs explicit boundary fixtures. In the inspected
zero-time helper, the first nonzero index is $i$ and the deletion interval is
$\texttt{start}:(\texttt{start}+i-1)$. For a conventional one-based vector
with times $(0,1,2)$, the abstract trace removes indices one and two, leaving only
time two. If all times are zero, the first-nonzero search has no result and the
deletion branch is not entered. The already-sorted event-vector constructor also
passes its vector into this mutating helper; its comment that the input is unaffected
does not settle aliasing. These are source-local concerns with inspected constructor
and export paths, not reported outcomes from a Julia execution. The safe proposal
route is to reject unsupported zero-time inputs until the intended rule and a native
regression fixture are agreed; silently adding a small epsilon would change the
clock and conceal the question.~\citep{exec24_survivalother_juliastats,exec24_survivalother_reference}

SurvivalAnalysis.jl supplies another architecture. Its pinned version 0.1.0 declares
Julia compatibility \texttt{\^{}1.2} and MIT licensing. Its response hierarchy contains
right-, left- and interval-censored forms, but the inspected documentation explicitly
says that only right-censoring methods are included. In particular,
\texttt{Surv(start,stop)} represents interval censoring. The parametric fitting
dispatch takes a right-censored response and admits Weibull or exponential baselines.
The Hessian inverse and standard-error calculation have exception paths producing
NaNs. An adapter must report those outcomes as unavailable uncertainty, not replace
them with zero standard errors. PH prediction distinguishes a linear predictor from
a transformed survival distribution, which is useful for a typed output contract.
The stated candidate role is an original Julia parametric worker, subject to baseline,
optimisation and uncertainty checks; its response classes do not justify a universal
censoring claim.~\citep{exec24_survivalother_juliaanalysis}

SurvivalModels.jl, at its pinned 0.1.6 declaration with Julia 1.11 compatibility and
MIT licensing, develops parametric general-hazard alternatives and Cox and
nonparametric components. The inspected CoxDefault implementation groups tied times
before subtracting $d\log R$ from the log likelihood: this is a Breslow-style tied
denominator. It should not be expected to agree numerically with the preceding
Efron-style path for every tied data set. With risk scores $(1,2,3)$ and the first
two subjects failing together, the numerator is two and the total risk score is six.
The Breslow partial-likelihood contribution is $2/6^2=1/18$, whereas fractional
removal gives $2/[6(6-3/2)]=2/27$. A parity fixture must name its tie convention
before classifying a difference as a failure.~\citep{exec24_survivalother_juliamodels,exec24_survivalother_reference}

Its Kaplan--Meier constructor forms event increments at stored event/censor times.
The vector survival prediction takes their cumulative product; scalar prediction
uses only times strictly less than the requested time. For one event at time two
among two accounts, with the other censored at four, the abstract boundary values
are $S(2-)=1$ and $S(2)=1/2$. The inspected scalar and stored-grid routes thus encode
different sides of the jump, despite documentation describing the usual less-than-or-equal
product. The proposed wrapper must explicitly distinguish left limits from survival
after the event and test scalar/vector agreement under the chosen meaning. A package
name or a successful fit cannot resolve this discrepancy. Delegated SurvivalBase and
SurvivalDistributions behaviour remains a separate inspection obligation.
\citep{exec24_survivalother_juliamodels,exec24_survivalother_reference}

JointSurvivalModels.jl offers a different reuse candidate: constructing distributions
from hazards linked to a longitudinal process. Its version 1.0.0 snapshot declares
Julia 1 compatibility and MIT licensing. The inspected structure combines a baseline
hazard with covariate and linked-process effects. This can help formulate a stipulated
joint model of servicing history and event time, but a model of association does not
identify a payment-holiday intervention. Numerical support defaults to a positive
lower endpoint $10^{-4}$ and upper endpoint $10^4$; hazard integration begins at
the configured lower endpoint. The sampling method accepts a random generator but
the inspected uniform draw calls the unqualified global random function. Integration
tolerance, tail support and generator routing must therefore be inspected and tested
before reproducibility or distributional accuracy is claimed. A conditional modelling
worker is plausible; a blanket port would import these unresolved semantics.
\citep{exec24_survivalother_juliajoint}

\subsection{Rust causal labels require method-level inspection}
The Cameron-Lyons survival repository declares version 2.0.0, edition 2024,
Rust 1.94 and MIT licensing at the pinned commit. These are declarations in the
snapshot, not evidence of a tested toolchain or supported release. Its public Rust
module path re-exports target-trial, marginal-structural-model and TMLE-named
functions. Python registration is delegated and is not established merely by a
Rust export. This explicit distinction matters because an integration proposal must
name the entrypoint whose behaviour is being assessed.~\citep{exec24_survivalother_rustcore}

In the inspected target-trial function, a treated clone initiating within the grace
period receives follow-up $\min(\widetilde T-t_A,\tau)$, while the control clone
uses a clock censored at treatment initiation. With eligibility zero, treatment at
month two, event at month ten and grace period three, the untruncated treated duration
is eight rather than ten months from eligibility. With maximum follow-up five, the
treated duration becomes five but its event flag retains the original event status.
The authored trace is $(5,1)$, while an origin-clock observation administratively
censored at five would be $(5,0)$. This does not establish runtime behaviour of
every wrapper; it establishes why a native admission fixture must bind origin,
grace-period strategy, censoring time and event status jointly.
\citep{exec24_survivalother_rustcore,exec24_survivalother_reference}

The inspected censoring-weight routine ignores its censoring-covariate arguments,
assigns protocol-censored clones a constant based on the trimming option and can
then normalise the weights to mean one. It is not a fitted conditional censoring
model in this path. If two hypothetical conditional observation probabilities are
$1/2$ and $1/4$, their inverse weights are two and four and their mean-normalised
weights are $2/3$ and $4/3$. Giving both records ten and normalising yields one
and one; normalisation cannot create the missing conditioning information.
Likewise, a longitudinal weight's numerator model has a meaning distinct from mean
normalisation. Two treatment probabilities $1/4$ and $1/2$ give an unstabilised
weight eight; numerator probabilities $1/2$ and $1/2$ give a stabilised weight two,
whereas mean normalisation of a single weight gives one.
\citep{exec24_survivalother_rustcore,exec24_survivalother_reference}

The output's name and unit require care. The inspected
\texttt{risk\_difference} field contains a difference between restricted means,
so its unit is time. The main calculation weights clones and uses one hundred
grid points; its subject-resampling bootstrap recreates clones without the inspected
censoring-weight call and uses fifty grid points. Fewer than two bootstrap results
trigger a fixed pair of 0.5 standard errors. The sequential-trial helper skips
small eligible sets and retains successful results from a conditional call, so a
delivery wrapper must preserve attempted, skipped, failed and returned trial counts.
These findings motivate rejection of direct causal admission for these entrypoints
at this snapshot, not rejection of every component in the repository.
\citep{exec24_survivalother_rustcore}

The inspected survival-TMLE function assigns outcome one to observations surviving
beyond a horizon, zero to prior observed events and one half to observations censored
before that horizon. A half is an assumption about an unresolved outcome. With one
known survivor, one known event and two unresolved records among four subjects, the
finite-cohort survival fraction can range from $1/4$ to $3/4$; imputing one half for
each unresolved record chooses $1/2$. It does not identify that value. The inspected
function fits a propensity and an outcome regression but contains no explicit
censoring-model adjustment for that assignment. Its restricted-mean integration starts
at the first supplied horizon, while its reported mean-time standard error is an
unweighted combination of pointwise errors without time widths or cross-horizon
covariances. Rescaling months to days must rescale a time-valued standard error;
the abstract expression without time widths does not do so. A method name cannot
substitute for these checks, and no conclusion here extends to uninspected censoring
or TMLE alternatives elsewhere in the repository.
\citep{exec24_survivalother_rustcore,exec24_survivalother_reference}

\subsection{An interval-data candidate and a justified refusal}
The small Rust turnbull snapshot declares version 0.1.0 and edition 2021. No licence
grant was located in the inspected Cargo declaration or complete root tree, so
adoption or redistribution cannot be presumed. Its response type includes interval
and entry information, but the inspected estimator constructs support only from
uncensored exact-event times and an infinity point; handling truncation and variance
is marked unfinished. With only finite intervals $(1,2]$ and $(2,3]$, this abstract
support construction supplies only infinity. Neither interval contains that point,
so both membership rows have zero sum. The normalisation step then has no positive
denominator. This is a reproducible support witness, not a native Rust test result.
The visible convergence loop also lacks an explicit iteration cap.
\citep{exec24_survivalother_rustturnbull,exec24_survivalother_reference}

The defensible proposal disposition is reference-only pending source, rights and
algorithmic repairs. A qualified interval-censoring worker would need an appropriate
support construction, a declared truncation likelihood, finite-iteration and
nonfinite-value handling, and an uncertainty contract. A zero variance placeholder
must not become a claim of certainty. An interval-only fixture is especially valuable
because ordinary exact-event examples can succeed without exercising the capability
named by the package.

\subsection{Adoption choices, public prototypes and full costs}
The preferred architecture preserves selected original engines behind typed,
versioned workers. JuliaStats is a candidate for narrowly specified right-censored
estimation after boundary qualification; SurvivalAnalysis for its declared parametric
response family; SurvivalModels for explicit tie and prediction contracts; and
JointSurvivalModels for a stipulated joint distribution after numerical and random
generator checks. Rust components can be reference or future worker candidates,
but the inspected causal entrypoints require semantic repair before admission.
The interval-data prototype remains reference-only while its licence and support
questions are unresolved. Whole-library ports are not justified by language
preference alone: a port transfers maintenance, rights, numerical and semantic
obligations without supplying identification.

A generated/public prototype can test the merit of this architecture before an
expensive bank experiment. Register exact-event, all-censored, zero-time, delayed-entry,
interval-only, tied-event, competing-event, administrative-censoring and malformed-input
fixtures; match each fixture to a target-specific oracle. Compare scalar and vector
outputs, time-unit rescaling, row permutations, random-seed routing and complete
estimator resampling. Separately evaluate uncertainty under a declared independent
or clustered sampling law. Finite-law examples establish those exact laws; public
data establish only the observable behaviours and assumptions they actually support.
Neither silently supplies counterfactual bank truth.

The CLI/MCP result should carry the source snapshot, environment and rights profile,
original row and subject identities, target and clock, tie/boundary options, diagnostics,
refusal reason, uncertainty method and complete costs. Preserve unsuccessful requests
and human repairs in the denominator. In a stipulated accounting example, twelve
initiated requests costing USD 36 and producing three accepted numeric deliveries
cost USD 3 per initiation but USD 12 per accepted delivery. Correct refusals should
be scored separately from accepted numeric results; deleting them would distort
both reliability and cost. These amounts are authored examples, not vendor prices
or measured performance.~\citep{exec24_survivalother_reference}

The six selected profiles constitute a contrastive review, not exhaustive
other-language coverage. The Julia discovery tail, Rust alternatives, failed
target-trial search and uninspected delegates remain explicit follow-ups. The method
owner must approve the target and observation law; the adapter owner must preserve
units, failures and identities; the rights owner must clear the exact distribution;
and the evaluation owner must preregister public/generated comparisons. This gives
concrete next actions while leaving native, exhaustive and bank acceptance unestablished.


The chapter's common requirement is to preserve the same account history
through every stage of the calculation. Changing the time origin, event
priority, policy law, averaging population or unit of uncertainty can change
the financial decision while leaving a plausible curve on the screen. The
mathematical references and source-specific failures show which distinctions
a proposed worker must expose. The next treatment compares regimes,
heterogeneous effects and continuous doses using those requirements, rather
than assuming that one longitudinal label makes the methods interchangeable.

